%%%BLOCK id=021-01 ch=12 sec=磁通量·ΔΦ分类 kind=knowledge alt=none cont=none y=0.03-0.18
\begin{center}
$\Phi=BS\cos\theta$\\
无匝数
\end{center}
\begin{itemize}
  \item $S$变　动生　$\Delta\Phi=\Delta S\cdot B\sin\theta$
  \item $B$变　感生　$\Delta\Phi=S\cdot\Delta B\sin\theta$
  \item $SB$变　动感　$\Delta\Phi=\Phi_2-\Phi_1$，$\ne\Delta S\,\Delta B\sin\theta$
  \item $\theta$变　正弦式交流电（平动切割）
% CHECK: 原稿写“平动切割”，θ变化对应线圈转动切割，疑为笔误；另 Φ=BScosθ 与下面各式 sinθ 中 θ 的含义不同，照原样
  \item 多磁场叠合　$\Delta\Phi=\Phi_2-\Phi_1$
\end{itemize}
%FIG p021_a 0.775 0.05 0.90 0.13
\notefig[0.25]{figures/p021_a.png}
%%%END
%%%BLOCK id=021-02 ch=12 sec=楞次定律·套住磁铁看抵消 kind=method alt=none cont=none y=0.18-0.40
\textbf{套住磁铁看抵消}

磁感线\unclear{条}，内$=$外
%FIG p021_b 0.07 0.235 0.29 0.344
\notefig[0.3]{figures/p021_b.png}
管外磁场抵消管内磁场

抵消的多\\
$\Phi$小
%%%END
%%%BLOCK id=021-03 ch=12 sec=动生电动势·单杆平动 kind=formula alt=none cont=none y=0.40-0.52
\[
\left\{\begin{aligned}
F_A&=BIL\\
I&=\frac{E}{R_{\text{总}}}\\
E&=BLv
\end{aligned}\right.
\qquad F_{\text{安}}=\frac{B^2L^2v}{R}\quad\text{动生单杆平动}
\]
%%%END
%%%BLOCK id=021-04 ch=12 sec=电磁感应·电荷量 kind=formula alt=none cont=none y=0.52-0.59
\[
q=\bar I\,t,\qquad I=\frac{\bar E}{R},\qquad \bar E=n\frac{\Delta\Phi}{\Delta t},\qquad q=n\frac{\Delta\Phi}{R_{\text{总}}}
\]
（$q=n\frac{\Delta\Phi}{R_{\text{总}}}$ 中的 $n$ 处有箭头 $\downarrow$ 指向“匝数”）
%%%END
%%%BLOCK id=021-05 ch=12 sec=电磁感应·能量 kind=formula alt=none cont=none y=0.59-0.67
\[
W_{A\text{负}}+W_{\text{电池放电}}=W_{A\text{正}}+Q
\]
% CHECK: 下标“负/正”为按字形推测（疑指安培力做负功/正功）
（“电池放电”项正下方有小注：）(无)\\
干电池
%%%END
%%%BLOCK id=021-06 ch=12 sec=两级线圈·电流变化 kind=knowledge alt=none cont=none y=0.67-0.78
\textbf{两级线圈}

大小：$I_1$ 求导 $I_2$　满足楞次定方向
% CHECK: 原稿“满足楞次定方向”，疑漏字（楞次定律判方向），照原样

瞬时值：
\begin{center}
\begin{tabular}{ll}
$\sin$ & \unclear{求导}$\;\cos$\\
$\cos$ & $-\sin$
\end{tabular}
\end{center}
%%%END
%%%BLOCK id=021-07 ch=12 sec=动生电动势·有效切割长度 kind=method alt=04 cont=none y=0.78-0.89
\textbf{动生}
%FIG p021_c 0.075 0.82 0.245 0.89
\notefig[0.3]{figures/p021_c.png}
\[
E=B\frac{l}{\sin\theta}\,v\sin\theta
\]
%FIG p021_d 0.42 0.775 0.578 0.93
\notefig[0.3]{figures/p021_d.png}
平抛\quad $E=BLv_0$
%%%END
%%%BLOCK id=021-08 ch=12 sec=动生电动势·等效电源 kind=diagram alt=11 cont=none y=0.89-0.97
% CHECK: 下面两幅图无文字说明，疑为含电源的动生回路（R、r）与电磁流量计/磁流体发电类装置（d）
%FIG p021_e 0.04 0.895 0.205 0.965
\notefig[0.3]{figures/p021_e.png}
%FIG p021_f 0.26 0.905 0.44 0.96
\notefig[0.3]{figures/p021_f.png}
%%%END
%%%BLOCK id=022-02 ch=12 sec=电源串并联·等效电源 kind=formula alt=10 cont=none y=0.14-0.30
\textbf{电源串联}
\begin{align*}
E_{\text{总}}&=E_1\pm E_2\quad\text{（看方向）}\\
r_{\text{总}}&=r_1+r_2
\end{align*}
\textbf{电源并联}
%FIG p022_b 0.20 0.21 0.32 0.295
\notefig[0.25]{figures/p022_b.png}
\begin{align*}
E_{\text{总}}&=\max\{E_1,E_2\}\\
r_{\text{总}}&=\frac{r_1r_2}{r_1+r_2}
\end{align*}
% CHECK: 原稿并联“E总”式中 max 写在花括号下方；物理上并联等效电动势一般并非 max{E1,E2}，照原样
%%%END
%%%BLOCK id=022-03 ch=12 sec=感生电动势·感应部分当电源 kind=knowledge alt=none cont=none y=0.31-0.42
\textbf{感生}　感应部分当电源，其余部分为外电路，电源两端电压为路端电压
\[
\varepsilon=n\frac{\Delta\Phi}{\Delta t}\qquad
\Delta\Phi=S\cdot\Delta B\sin\theta\ \Leftarrow\ n\frac{\Delta B}{\Delta t}S\sin\theta\qquad
\text{BS夹角}
\]
（$S$ 下方有箭头 $\uparrow$ 标注：）磁场和线圈的重合面积
%%%END
%%%BLOCK id=022-04 ch=12 sec=感生电场·涡旋电场 kind=knowledge alt=none cont=none y=0.43-0.57
\textbf{感生电场（涡旋电场）}

$B$变 $\to$ $E_{\text{感}}$ 不是 静电场
\begin{itemize}
  \item 电场线：闭合同心圆
  \item $E$方向：和 $I_{\text{感}}$ 一样
  \item 做功　$W_{\text{电}}=q\varepsilon=q\dfrac{\Delta\Phi}{\Delta t}$
  \item $E$大小　$qE\cdot2\pi r=W_{\text{电}}=q\varepsilon=q\dfrac{\Delta\Phi}{\Delta t}$
\end{itemize}
通过做功算　求出$E$
%%%END
%%%BLOCK id=022-05 ch=12 sec=感感·两区域感生叠加 kind=method alt=none cont=none y=0.57-0.66
\textbf{感感}
%FIG p022_c 0.17 0.565 0.325 0.65
\notefig[0.3]{figures/p022_c.png}
\begin{align*}
E&=E_1+E_2\\
E&=n\frac{\Delta B_1}{\Delta t}S_1\sin\theta+n\frac{\Delta B_2}{\Delta t}S_2\sin\theta
\end{align*}
%%%END
%%%BLOCK id=022-06 ch=12 sec=感动·感生与动生叠加 kind=method alt=none cont=none y=0.66-0.99
\textbf{感动}
\begin{itemize}
  \item 切恒定$B$
%FIG p022_d 0.31 0.672 0.445 0.728
\notefig[0.3]{figures/p022_d.png}
\[
E=|E_1-E_2|\qquad\text{·增顺　}\times\text{增逆}
\]
\[
E=\left|n\frac{\Delta B_1}{\Delta t}S_1\sin\theta-B_2Lv_0\right|
\]
  \item 切变化$B$
%FIG p022_e 0.195 0.735 0.335 0.815
\notefig[0.3]{figures/p022_e.png}
%FIG p022_f 0.385 0.738 0.565 0.804
\notefig[0.3]{figures/p022_f.png}
\end{itemize}
假定$B$不变　假设棒不动
\[
E_1=BLv_0\qquad E_2=n\frac{\Delta B}{\Delta t}S\sin\theta
\]
求 $t_1$ 时刻 $E_{\text{电动势}}$
\begin{align*}
E_1&=BLv_0=kt_1Lv_0\\
E_2&=1\cdot k\cdot Lv_0t_1\cdot1
\end{align*}
% CHECK: 原稿 E2=1·k·LV0t1·1，两端的“1”可能是 n=1 与 sinθ=1，也可能是绝对值竖线，照原样按数字1转录
\[
E=E_1\pm E_2
\]
%%%END
%%%BLOCK id=023-01 ch=12 sec=单杆模型·分析方法 kind=method alt=none cont=none y=0.075-0.20
\textbf{单杆}　过程/状态/动能
\[
\left\{\begin{aligned}
F_A\ \text{恒力}\quad&W_A=\pm BILx\\
F_A\ \text{变力}\quad&\pm W_A
\end{aligned}\right.
\]
冲量：
\begin{align*}
I_A&=BIL\cdot\Delta t\quad\text{过程写动\unclear{能}动量}\\
I_A&=\pm BLq
\end{align*}
% CHECK: “过程写动能/动量”处“能”与“量”字迹重叠，疑为把“动能”改成“动量”（动量定理），照原样
状态写牛二　最大　最小　稳定$\Leftrightarrow a=0$
%%%END
%%%BLOCK id=023-02 ch=12 sec=电磁感应·能量 kind=knowledge alt=none cont=none y=0.20-0.31
\begin{itemize}
  \item 功和功关系　动能定理
  \item 功和能关系　动能定理
% CHECK: 原稿第二行右侧同为“动能定理”（“功”字有蓝色荧光笔涂抹），疑应为“功能关系”，照原样
  \item 能和能关系　能量守恒
  \begin{itemize}
    \item 机械能　机械能定理
  \end{itemize}
\end{itemize}
%%%END
%%%BLOCK id=023-03 ch=12 sec=电磁感应·电热分配 kind=method alt=10 cont=none y=0.31-0.62
\textbf{电热/功/功率　分配}
%FIG p023_a 0.075 0.365 0.215 0.442
\notefig[0.3]{figures/p023_a.png}
%FIG p023_b 0.23 0.362 0.35 0.435
\notefig[0.3]{figures/p023_b.png}
\begin{center}
\begin{tabular}{p{0.44\linewidth}p{0.44\linewidth}}
并联　$Q$与$R$成反比 & 串联　$Q$与$R$成正比\\[2pt]
$\dfrac{U^2}{R}t$ & $I^2Rt$\\[6pt]
$Q_1=\dfrac{R_2}{R_1+R_2}Q$ & $Q_1=\dfrac{R_1}{R_1+R_2}Q$\\[8pt]
$Q_2=\dfrac{R_1}{R_1+R_2}Q$ & $Q_2=Q\dfrac{R_2}{R_1+R_2}$\\[8pt]
$Q$和其它电阻成正比 & $Q$和自身电阻成正比
\end{tabular}
\end{center}
$P$、$W$ 一样
%%%END
%%%BLOCK id=023-04 ch=12 sec=电磁感应·变力冲量(微积分) kind=derivation alt=90 cont=none y=0.62-0.72
\textbf{微积分}
\[
\left\{\begin{aligned}
&f=kv\\
&qvB\\
&\frac{B^2L^2v}{R}
\end{aligned}\right.
\qquad a\text{变化的运动}
\]
% CHECK: “a变化的运动”首字形似 Q/a，按“a（加速度）变化的运动”转录
\[
I=\int f\,\dd t=k\,v\,\dd t=kx
\]
% CHECK: 原稿第二个等号后为“k vdt”，无积分号，照原样
%%%END
%%%BLOCK id=023-05 ch=12 sec=电磁感应·动量 kind=derivation alt=none cont=none y=0.72-0.99
\[
I_A=\pm BLq=\pm\frac{B^2L^2x}{R}=m(v_2-v_1)
\]
（$BLq$ 下方标注：）流过导体电荷量\\
（箭头 $\uparrow$ 指向 $x$，标注：）单杆导棒位移/双杆位移差
\begin{align*}
-BIL\,\Delta t&=m\,\Delta v\\
-BL\,\Delta q&=m\,\Delta v\\
-\frac{B^2L^2v}{R_{\text{总}}}\,\Delta t&=m\,\Delta v\\
-\frac{B^2L^2x}{R_{\text{总}}}&=m\,\Delta v
\end{align*}
求和：
\begin{align*}
-BLq&=m(v_2-v_1)\\
-\frac{B^2L^2x}{R_{\text{总}}}&=m(v_2-v_1)
\end{align*}
%%%END
%%%BLOCK id=024-01 ch=12 sec=单杆模型·动源 kind=method alt=10 cont=none y=0.00-0.10
\textbf{动源}
%FIG p024_a 0.056 0.0 0.275 0.088
\notefig[0.3]{figures/p024_a.png}
\[
I=\frac{E-BLv}{R+r}\qquad\text{最终}\ E=BLv\quad I=0\quad\text{导棒匀速}\quad v=\frac{E}{BL}
\]
%%%END
%%%BLOCK id=024-02 ch=12 sec=单杆模型·动容 kind=method alt=09 cont=none y=0.11-0.32
\textbf{动容}　\textbf{恒力充电型}
%FIG p024_b 0.025 0.146 0.215 0.232
\notefig[0.3]{figures/p024_b.png}
\begin{align*}
F-BIL&=ma\\
F-B\frac{\Delta Q}{\Delta t}L&=ma\\
\Delta Q&=C\,\Delta U=C\,\Delta(BLv)=CBL\,\Delta v\\
F-BCBL^{2}a&=ma\\
a&=\frac{F}{m+CB^2L^2}\ \text{定值}\quad\text{匀加速直线运动}
\end{align*}
% 原稿“匀加速直线运动”下有波浪线下划线
%%%END
%%%BLOCK id=024-03 ch=12 sec=单杆模型·动容 kind=method alt=09 cont=none y=0.32-0.49
\textbf{初速充电型}
%FIG p024_c 0.07 0.325 0.315 0.462
\notefig[0.4]{figures/p024_c.png}
最终 $I=0$　导棒匀速　流过导棒电荷量 $\Delta Q=Q$
\[
\left\{\begin{aligned}
-BLQ&=mv-mv_0\\
\frac{Q}{C}=U&=BLv
\end{aligned}\right.
\quad\Rightarrow\quad Q,\ v,\ U
\]
\[
mv_0=(m+CB^2L^2)v\qquad v=\frac{mv_0}{m+CB^2L^2}
\]
% CHECK: 原稿“mv0=(m+CB²L²)v”右括号缺失，按其意补全
%%%END
%%%BLOCK id=024-04 ch=12 sec=单杆模型·动容 kind=method alt=09 cont=none y=0.49-0.65
\textbf{含容放电}
%FIG p024_d 0.015 0.508 0.345 0.64
\notefig[0.4]{figures/p024_d.png}
\[
\overset{\downarrow}{I}=\frac{\overset{\downarrow}{U}-\overset{\uparrow}{BLv}}{R}
\]
最后　$I=0$ 后　导棒匀速　$F_A=0$
\[
\left\{\begin{aligned}
+BL(Q_0-Q)&=mv-0\\
\frac{Q}{C}=U&=BLv
\end{aligned}\right.
\quad\Rightarrow\quad Q,\ v,\ U
\]
%%%END
%%%BLOCK id=024-05 ch=12 sec=单杆模型·动容 kind=method alt=09 cont=none y=0.65-0.68
不充电不放电型　电容器为理想电压表　$Q=CU$
%%%END
%%%BLOCK id=024-06 ch=12 sec=涡流·动生涡流 kind=method alt=none cont=none y=0.69-0.93
\textbf{动生涡流}
\begin{itemize}
  \item 有$B$位置　相当于电源
  \item 无$B$位置　相当于外电路
\end{itemize}
%FIG p024_e 0.09 0.765 0.225 0.845
\notefig[0.25]{figures/p024_e.png}
实心圆盘

有$E$　无$I$　无$Q$

无电磁阻尼 $\to$ 不减速(匀速)
%FIG p024_f 0.375 0.76 0.515 0.86
\notefig[0.25]{figures/p024_f.png}
有$E$，有$I$　有$Q$

有电磁阻尼 $\to$ 减速
%FIG p024_g 0.605 0.765 0.76 0.836
\notefig[0.3]{figures/p024_g.png}
$\downarrow$ 不能使灯泡亮，电流流不出来

有$E$　$I$　$Q$　涡流

有电磁阻尼 $\to$ 减速
%%%END
%%%BLOCK id=024-07 ch=12 sec=涡流·感生涡流 kind=knowledge alt=none cont=none y=0.92-1.00
% 页面底部被截断，立方体图不完整
\textbf{感生涡流}

无\unclear{数}层、闭合\unclear{电路}

类似滑动摩擦力
%FIG p024_h 0.185 0.929 0.31 1.0
\notefig[0.25]{figures/p024_h.png}
%%%END
%%%BLOCK id=024-08 ch=12 sec=涡流·利用与防止 kind=knowledge alt=none cont=none y=0.92-1.00
利用涡流
\begin{itemize}
  \item 热效应：电磁炉(金属锅、真空冶炼钢)
  \item 磁效应：金属探测仪、探伤、扫雷
\end{itemize}
防止涡流：变压器铁芯：采用电阻率大的硅钢片
%%%END
%%%BLOCK id=025-01 ch=12 sec=涡流·电磁阻尼 kind=knowledge alt=none cont=none y=0.06-0.13
\textbf{电磁阻尼}

导体动，磁场动　都有阻碍相对运动现象
%%%END
%%%BLOCK id=025-02 ch=12 sec=涡流·电磁驱动 kind=knowledge alt=none cont=none y=0.14-0.21
\textbf{电磁驱动}

磁体/导体　一个先转，产生涡流　从而带动另一个转动（滞后、后慢）
%%%END
%%%BLOCK id=032-02 ch=12 sec=自感·含感电路 kind=knowledge alt=13,10 cont=none y=0.10-0.36
\textbf{含感电路}

\noindent 自感电动势\quad $E=L\dfrac{\Delta I}{\Delta t}$\quad $L$ 与匝数$\uparrow$、粗细$\uparrow$、长度$\uparrow$、铁芯有$\uparrow$无$\downarrow$ 有关\quad 方向：阻碍流过自身电流的变化% 原稿中箭头写在“匝数”“粗细”“长度”“有无”上方

\noindent\hspace*{4em}自感电流方向与原电流：增反减同

\noindent 在直流电路：自感在
\begin{itemize}
\item 通电瞬间：$R$ 从 $\infty\to0$\textsuperscript{动态电阻}\quad 注\ 串反并同/串反并\unclear{恒}
\item 电流稳态：$R_L=0$\textsuperscript{导线}\ or\ $R_L$ 很小$\neq0$
\item 断电瞬间：$E$ 减小\textsuperscript{动态电源}\quad $I_{\text{自感}}$ \unclear{$\le$} $I_{\text{原}}$
\end{itemize}

\noindent 自感在交流电阻中\ 视为电阻（感抗）\ $R_L=2\pi Lf$\quad 通直阻交% CHECK: “交流电阻中”疑为“交流电路中”(原稿如此)
%%%END
%%%BLOCK id=052-01 ch=12 sec=楞次定律·右手定则 kind=summary alt=none cont=none y=0.03-0.58
\textbf{电磁感应}

\begin{itemize}
\item 电磁感应现象、楞次定律
 \begin{itemize}
 \item 磁通量：\ $\phi=\vec B\cdot\vec S=B\cdot\underline{S\cos\theta}$，\ $\uparrow$ 垂直于磁场的面积
 \item 电磁感应现象：\ $\phi$变 $\to U_{\text{感}}$ + 闭合回路 $\to I_{\text{感}}$
 \item 感应电流方向的判定，\ 楞次定律、右手定则、
 \end{itemize}
\end{itemize}

电磁感应现象的理解

楞次定律和右手定则的应用 \quad 增反减同\ 来拒去留

楞次定律的推论的应用 \quad 增缩减扩\ 增反减同、

三定则一定律的综合应用 \quad 安培定则、左手定则、右手定则\ 楞次定则.
% 原稿此项右侧分两行写：第一行「安培定则、左手定则、」，第二行「右手定则 楞次定则.」（第二行与「二次感应问题」同在一行）
% CHECK: 原稿写作“楞次定则”，疑为“楞次定律”笔误

二次感应问题 \quad 程序法、倒推法
%%%END
%%%BLOCK id=052-02 ch=12 sec=法拉第电磁感应定律 kind=summary alt=none cont=none y=0.57-0.72
\begin{itemize}
\item 法拉第电磁感应现象\ 自感、涡流.
 \begin{itemize}
 \item 法拉第电磁感应定律 \quad $E=N\dfrac{\Delta\phi}{\Delta t}$
 \item 自感和\underline{涡流}. \quad $E=NL\dfrac{\Delta I}{\Delta t}$ \quad $L$与线圈大小、形状、圈数\ 铁芯有无\ 有关\quad 单位亨利$H$
 % CHECK: 自感电动势一般为 E=LΔI/Δt，原稿看起来写成 E=NL ΔI/Δt（多一个 N）
 \begin{itemize}
 \item $\uparrow$ 非匀强磁场.
 \item 锅自身发热，而非加热源，只\unclear{能}用铁锅
 \item 旋涡状电流 流不出去\ 不能点亮灯泡
 \end{itemize}
 \end{itemize}
\end{itemize}
%%%END
%%%BLOCK id=052-03 ch=12 sec=法拉第电磁感应定律 kind=formula alt=none cont=none y=0.69-0.95
\textbf{法拉第电磁感应定律的理解运用.} $\swarrow$
% 原稿此处有一箭头(↙)指向下方左侧各组公式

S变：
\[
E=n\frac{B\Delta S}{\Delta t}=BLv
\]
B变：
\[
E=n\frac{\Delta B}{\Delta t}S=nKS
\]
BS一起变：
\[
E=n\frac{|B_2S_2-B_1S_1|}{\Delta t}
\]
\[
\ne n\frac{\Delta B\,\Delta S}{\Delta t}
\]
\[
q=\bar I\,\Delta t=\frac{\bar E}{R}\Delta t=n\frac{\Delta\phi}{\Delta t\,R}\Delta t
\]
\[
=n\frac{\Delta\phi}{R}
\]
%%%END
%%%BLOCK id=052-04 ch=12 sec=动生电动势·BLv kind=method alt=none cont=none y=0.69-0.90
\textbf{$E=BLv$的理解及运用}\quad（正交性 $BLv$两两垂直、有效性 $L$、相对性 $v$）

%FIG p052_a 0.65 0.715 0.772 0.763
\notefig[0.25]{figures/p052_a.png}

$B\perp v$ \quad $L$有效长度 \quad $v_{\text{中}}$、$v$为相对磁场的速度

\[
v_{\text{中}}=AC\,\text{\unclear{上}}\qquad E=B\cdot AC\cdot\frac{v_A+v_C}{2}
\]
\[
=B\cdot\overline{AC}\cdot\frac{\omega\,\overline{OA}+\omega\,\overline{OC}}{2}
\]
% 第一式 v_中 处有涂改（“=”压在 v 与 中 之上）
%%%END
%%%BLOCK id=052-05 ch=12 sec=自感·涡流·电磁阻尼与驱动 kind=summary alt=none cont=none y=0.80-0.92
自感电动势 $\swarrow$ 阻碍原电流变化\quad （延缓）.\ \unclear{原}电流$I$大\ 并联灯闪亮
% “并”字被涂抹覆盖，其余按原稿读出

自感、涡流，电磁阻尼与电磁驱动，

\begin{tabular}{@{}lll@{}}
 & 导体动 & 磁场动（棒的滞后性）\\
感应电流 & 安培力 阻碍 & 安培力 动力\\
 & 动能$\to$电能$\to$热能 & 磁场能$\to$电能$\to$机械能
\end{tabular}
%%%END
%%%BLOCK id=052-06 ch=12 sec=电路图像·动力学动量能量 kind=summary alt=none cont=none y=0.91-1.00
电磁感应中的电路与图像问题 \quad 排除法\ 函数法.
\[
I=\frac{E}{R+r}\quad U=\frac{RE}{R+r}\quad P=IU\quad Q=I^2Rt\quad q=CU\quad E=BLv=n\frac{\Delta\phi}{\Delta t}=\frac12BL\omega^2=BLv_{\text{\unclear{平}}}
\]
% 末项 BLv 右下角的小标记不清，暂读作“平”

电磁感应的动力学、动量、能量问题、三把金钥匙 \quad $q=n\dfrac{\Delta\phi}{R}$
%%%END
%%%BLOCK id=054-01 ch=12 sec=电磁驱动·磁悬浮模型 kind=derivation alt=none cont=none y=0.03-0.37
% 手写 v 与 V 不分，这里统一记为小写 v
\textbf{电磁驱动，带动滞后}

\unclear{$I_f=kx$} $\downarrow$
% 页面右上角（No.栏右侧）的小字，箭头向下指向右侧的 f=kv

\textbf{磁悬浮}（二倍电流\ 四倍力）

%FIG p054_a 0.0 0.068 0.283 0.137
\notefig[0.3]{figures/p054_a.png}

磁场速度$v_B$，安培力驱动

%FIG p054_b 0.545 0.058 0.685 0.14
\notefig[0.25]{figures/p054_b.png}

$a\downarrow$到0后匀速直线运动，$f=kv$ \quad \red{框有$n$匝}

开始：
\[
E=\red{n}\,2\cdot BLv\qquad I=\red{n}\,\frac{E}{R}
\]
\[
F_{\text{安}}=\red{n}BI\cdot2L=\frac{\red{n^2}\,4B^2L^2v_B}{R}
\]
启动后：
\[
E=\red{n}\,2BL(v_B-v_{\text{框}})
\]
\[
F_{\text{安}}=\frac{\red{n^2}\,4B^2L^2(v_B-v_{\text{框}})}{R}
\]
由牛顿第二定律\ $F_{\text{安}}-f=ma$\quad $a=0$时速度最大.
\[
v_m=\frac{4B^2L^2v_B\cdot\red{n^2}}{KR+4B^2L^2}
\]
每秒的$E_{\text{补\unclear{充}}}=Kv_m\cdot v_m+I^2R$，\quad $I=\dfrac{\red{n}\,2BL(v_B-v_m)}{R}$

（功率）\qquad 能量 由磁场提供

\[
E=\red{n}BLv_{\text{相}}
\]
\struck{$q=nBLv$}
\[
I_{\text{冲}}=nBLq=\frac{n^2B^2L^2x}{R}
\]
%%%END
%%%BLOCK id=054-02 ch=12 sec=电磁驱动·磁悬浮模型 kind=problem alt=none cont=none y=0.37-0.71
\begin{problem}
框受阻力$f$，磁场由静止开始匀加速$t_1$后，框\unclear{……}正在向右匀加速，速度为$v_1$，求框何时运动？
% 原稿“框”字之后有一小段被涂改液盖住，字迹不可辨
\begin{solution}
对框\quad
$\left\{\begin{aligned}
&F_{\text{安}}-f=ma\\
&v_B=at_1\\
&v_{\text{框}}=v_1
\end{aligned}\right.$

$a_1=a_2$.
\[
F_{\text{安}}=\frac{4B^2L^2(at_1-v_1)}{R}
\]
\[
\frac{4B^2L^2(at_1-v_1)}{R}-f=ma
\]
\[
F_{\text{安}}=\frac{4N^2B^2L^2v_{\text{相}}}{R}
\]
\[
v_{\text{相}}=\bigl(v_B+a\Delta t-(v_{\text{框}}+a\Delta t)\bigr)\ \text{不变.}
\]
刚开始运动时
\[
f=\frac{4N^2B^2L^2at_0}{R}
\]
故
\[
t_0=\frac{fR\,(4N^2B^2L^2-mR)}{4N^2B^2L^2\,(fR+4N^2B^2L^2v_1)}
\]

\fcolorbox{red!70!black}{white}{\parbox{0.9\linewidth}{
\red{演草}\quad（原稿用红笔框住此区域）
\[
\text{\red{\textcircled{a}}}=\frac{4B^2L^2(\text{\red{\textcircled{a}}}t_1-v_1)-fR}{mR}\qquad \text{\red{没解完}}
\]
开始移动时\quad $F_0-f=0$\qquad $F_0=\dfrac{4B^2L^2at_0}{R}$
\[
\frac{4B^2L^2at_0}{R}=f
\]
\[
t_0=\frac{fR^2m}{4B^2L^2\bigl(4B^2L^2(at_1-v_1)-fR\bigr)}
\]
\[
a=\frac{fR+4B^2L^2v_1}{4B^2L^2t_1-mR}
\]
}}
\end{solution}
\end{problem}
%%%END
%%%BLOCK id=054-03 ch=12 sec=电磁驱动·磁悬浮模型 kind=problem alt=none cont=none y=0.67-0.98
\begin{problem}
达到$v_m$后，磁场突然停止，$t$后框也停止，求框在$t$内移动距离.
\begin{solution}
\[
\sum F\cdot\Delta t+ft_1=mv_m
\]
% CHECK: 原稿写作 ft_1，按题意（磁场停止后 t 时间框停止）此处时间似应为 t
\[
\frac{4B^2L^2x}{R}+ft_1=mv_m
\]
\[
I_{\text{安}}=\frac{n^2B^2L^2x}{R}
\]
\[
nBLq,\qquad \frac{n^2B^2L^2x}{R}
\]
\end{solution}
\end{problem}
%%%END
%%%BLOCK id=054-04 ch=12 sec=电磁感应·三大观点 kind=summary alt=none cont=none y=0.71-1.00
\textbf{三大观点}（原稿框内）

%FIG p054_c 0.485 0.715 1.0 1.0
\notefig[0.7]{figures/p054_c.png}

\begin{description}
\item[动力学]
\begin{tabular}{@{}ll@{}}
源 & $E=n\dfrac{\Delta\Phi}{\Delta t}$ \unclear{或} $E=BLv$\\
路 & 内电路\ 串联\quad 外电路\ 并联\quad $E=IR+Ir$\\
力 & $F_{\text{安}}=BIL\ \to\ F_{\text{合}}=ma$\\
运动 & $a$、$v$、$t$
\end{tabular}
% 原稿左侧有大括号括住 源/路/力/运动；箭头：E=BLv ← a、v（大弧线）；E=IR+Ir → F安=BIL；F合=ma → a、v、t
\item[动量]
$q=\bar It$

$B\bar IL\Delta t=mv_2-mv_1$

$BLq=mv_2-mv_1$

$Ft=I_{\text{冲}}$ $=mv_2-mv_1$
\[
-B\bar IL\Delta t=\frac{-B^2L^2\bar v\,\Delta t}{R_{\text{总}}}=0-mv_0.
\]
\[
\frac{-B^2L^2x}{R_{\text{总}}}=0-mv_0.
\]
求 $q$、$v$、$t$、$x$.
\item[能量]
其他能$\to$电能$\to$热能.\quad $Q=I^2Rt$ 或 $Q=W_{\text{克服安培力}}$ 或 $Q=\Delta E_{\text{其他能减少量}}$
\end{description}
%%%END
%%%BLOCK id=055-01 ch=12 sec=双棒模型 kind=method alt=none cont=none y=0.03-0.25
\textbf{双杆}

\begin{itemize}
\item $V_0$可由其他能量转化而来
\item 等宽\ 安培力抵消、动量守恒
\item 双向磁场\ 不抵消、动量不守恒
\end{itemize}
% 后两条写在页面右上角

① \textbf{等宽初速}

%FIG p055_a 0.122 0.088 0.268 0.238
\notefig[0.3]{figures/p055_a.png}

由动量守恒、能量守恒
\[
m_1V_{01}+m_2V_{02}=(m_1+m_2)V_{\text{共}}
\]
\[
Q_{\text{电}}+\frac12(m_1+m_2)V_{\text{共}}^2=\frac12m_1V_{01}^2+\frac12m_2V_{02}^2
\]

%FIG p055_b 0.29 0.166 0.63 0.24
\notefig[0.6]{figures/p055_b.png}
%%%END
%%%BLOCK id=055-02 ch=12 sec=双棒模型 kind=method alt=none cont=none y=0.24-0.51
② \textbf{不等宽初速}\quad \textbar 电源平衡\textbar\ 单体动量定理\ \textbar 能量守恒\textbar
% “单体”二字写在“动量定理”上方

%FIG p055_c 0.09 0.25 0.222 0.37
\notefig[0.25]{figures/p055_c.png}

\[
\left\{\begin{aligned}
&BL_1V_1=BL_2V_2 && \text{①}\\
&-BL_1q=m_1V_1-m_1V_{01} && \text{②}\\
&BL_2q=m_2V_2-m_2V_{02} && \text{③}\\
&\tfrac12m_1(V_{01}^2-V_1^2)+\tfrac12m_2(V_{02}^2-V_2^2)=Q_{\text{电}}
\end{aligned}\right.
\]
% 原稿第四式首项 m_1 写作大写 M_1
$\dfrac{\text{②}}{\text{③}}\quad -\dfrac{L_1}{L_2}=\dfrac{m_1(V_1-V_{01})}{m_2(V_2-V_{02})}$\ 得$V_1$、$V_2$

%FIG p055_d 0.17 0.415 0.36 0.508
\notefig[0.3]{figures/p055_d.png}
%%%END
%%%BLOCK id=055-03 ch=12 sec=双棒模型 kind=method alt=none cont=none y=0.50-0.70
\textbf{等宽恒力}\quad ③

%FIG p055_e 0.098 0.505 0.2425 0.695
\notefig[0.3]{figures/p055_e.png}

\[
\left\{\begin{aligned}
&F=(m_1+m_2)a\\
&Ft=m_1V_1+m_2V_2\qquad R=r_1+r_2\\
&\frac{B^2L^2(V_2-V_1)}{R}=m_1a\ \ \text{隔离1}\ \ \text{或}\ \ F-\frac{B^2L^2(V_2-V_1)}{R}=m_2a
\end{aligned}\right.
\]
% 原稿第一式 (m_1+m_2) 的 m 写作大写 M

$a_1=a_2$后\ $\Delta V$不再改变
\[
\Delta V=\frac{F_{\text{安}}(r_1+r_2)}{B^2L^2}\qquad a=\frac{F_{\text{安}}}{m}
\]

%FIG p055_f 0.24 0.635 0.42 0.72
\notefig[0.3]{figures/p055_f.png}
%%%END
%%%BLOCK id=055-04 ch=12 sec=双棒模型 kind=method alt=none cont=none y=0.69-0.90
\textbf{恒力摩擦型}\quad ④

%FIG p055_g 0.085 0.695 0.254 0.878
\notefig[0.3]{figures/p055_g.png}

$f_1=f_2=f$

若 $f<F<2f$\quad 2先变加速后匀速，1一直静止.

若 $F>2f$\quad 2先变加速，1后变加速，最终匀速.

%FIG p055_h 0.30 0.82 0.44 0.895
\notefig[0.25]{figures/p055_h.png}
%%%END
%%%BLOCK id=055-05 ch=12 sec=双棒模型 kind=method alt=none cont=none y=0.87-1.00
⑤ \textbf{不等宽双棒}\quad 恒力

%FIG p055_i 0.19 0.869 0.322 0.986
\notefig[0.3]{figures/p055_i.png}

电源平衡
\[
e=BL_2V_2-BL_1V_1\qquad \text{对}\,e\,\text{\unclear{求导}}\quad 0=BL_2a_2-BL_1a_1
\]
\[
\frac{a_1}{a_2}=\frac{L_2}{L_1}
\]
对2：\ $F-B\dfrac{e}{r_1+r_2}L_2=m_2a_2$

对1：\ $B\dfrac{e}{r_1+r_2}L_1=m_1a_1$
\[
\therefore F=\left(\frac{L_2^2}{L_1^2}m_1+m_2\right)a_2=\left(\frac{L_2}{L_1}m_1+m_2\frac{L_1}{L_2}\right)a_1
\]

%FIG p055_j 0.84 0.86 1.0 0.985
\notefig[0.25]{figures/p055_j.png}
%%%END
%%%BLOCK id=056-01 ch=12 sec=双杆模型·电荷量 kind=knowledge alt=none cont=none y=0.07-0.12
双杆中 $q$ 中 $x$ 为相对位移\quad 不等宽双杆中 $q=\dfrac{B\Delta S}{R}$，为相对面积 % 原稿“为相对面积”前未写主语(应指 ΔS)
%%%END
%%%BLOCK id=056-02 ch=12 sec=双杆模型·分热 kind=method alt=none cont=none y=0.12-0.33
\textbf{分热}
%FIG p056_a 0.10 0.118 0.24 0.19
\notefig[0.25]{figures/p056_a.png}
\[ Q_1=\frac{R}{R+R}\,Q_{\text{总}} \]
%FIG p056_b 0.09 0.19 0.292 0.315
\notefig[0.3]{figures/p056_b.png}
并联电路当整体来分配 % 图中下方两个圈：并联部分旁圈住 R/2(圈内左侧有小标记)，串联的 4Q 电阻下圈住 R
%%%END
%%%BLOCK id=056-03 ch=12 sec=双杆模型·分电荷 kind=method alt=none cont=none y=0.12-0.30
\textbf{分电荷}
%FIG p056_c 0.515 0.158 0.64 0.2295
\notefig[0.25]{figures/p056_c.png}
\[ q_1=q_2 \]
%FIG p056_d 0.50 0.2255 0.645 0.285
\notefig[0.25]{figures/p056_d.png}
\[ \frac{q_1}{q_2}=\frac{R_2}{R_1} \]
%%%END
%%%BLOCK id=056-04 ch=12 sec=双杆模型·等效电压 kind=method alt=none cont=none y=0.12-0.30
\textbf{等效电压}\quad 分压
%FIG p056_e 0.848 0.118 1.0 0.196
\notefig[0.28]{figures/p056_e.png}
\[ R_{\text{总}}=\frac{R}{2}+R \]
\[ V_{ab}=\frac{\dfrac{R}{2}}{R+\dfrac{R}{2}}\,BLv \]
%%%END
%%%BLOCK id=056-05 ch=12 sec=双杆模型·不等宽导轨 kind=problem alt=none cont=none y=0.34-0.69
% 原稿标题“非等宽受力棒 + 特殊条件”用方括号框起
\begin{problem}[非等宽受力棒 + 特殊条件]
%FIG p056_f 0.02 0.368 0.19 0.543
\notefig[0.25]{figures/p056_f.png}
① $t=0$ 开始用 $F$ 拉 $cd$，使其从静止开始沿导轨以 $a$ 匀加速，求 $cd$ 运动多久后 $ab$ 开始运动。
\begin{solution}
受力平衡 % 原稿“受力平衡”写在图右侧、公式左上方
\[ \mu mg=BIL_1=\frac{B\,BL_2at\,L_1}{2R}\ \to\ t \]
\end{solution}
② 用未知力 $F_0$ 拉 $cd$ 向右到 $V$ 后匀速，$ab$ 也恰好匀速，求 $V_{ab}$　$V_{cd}$　$P_{F_0}$。
\begin{solution}
受力 $\left\{\begin{array}{l}\text{对 2：}F_0=\mu mg+BIL_2\\ \text{对 1：}BIL_1=\mu mg\end{array}\right.$ \qquad $P_{F_0}=F_0V_2$

非电源\unclear{电动势}\quad $\varepsilon=BL_1V_1-BL_2V_2\neq 0$ % CHECK: “非电源”后两字辨认不清，可能是“电动势”
\end{solution}
③ $ab$ 运动到 $PQ$ 时刚好碰到\unclear{障碍}物停止，并将 $cd$ 上力改为 $F_1$，此时 $cd$ 为 $V_0$。

$L_2=\dfrac{1}{2}L_1$\qquad 一段时间后停止，求 $ab$ 停后 $cd$ 运动距离
\begin{solution}
由动量定理 \quad $\dfrac{-B^2L_2x_2}{\frac{3}{2}R}=0-mV_0$

电阻突变
%FIG p056_g 0.708 0.598 0.84 0.682
\notefig[0.25]{figures/p056_g.png}
\end{solution}
\end{problem}
%%%END
%%%BLOCK id=057-01 ch=12 sec=单杆·导轨电阻随位移变化 kind=problem alt=none cont=none y=0.02-0.55
\begin{problem}
%FIG p057_a 0.09 0.022 0.445 0.161
\notefig[0.4]{figures/p057_a.png}
$ac$、$bd$ 单位长度电阻 $r_0$，棒受 $F$ 向右匀加速，加速度为 $a$，到 $cd$ 时撤去 $F$，于 $ef$ 静止。
\begin{solution}
\[ E=BL\sqrt{2ax} \]
\[ mV-m\sqrt{2ax_1}=\frac{-B^2L^2x}{R+2r_0x_1} \]
\[ V=\sqrt{2ax_1}-\frac{B^2L^2x}{m\,(R+2r_0x_1)}\qquad \text{到 }x_2\text{ 时 }V_{\text{末}}=0 \]
\[ x_2=\frac{\sqrt{2ax_1}\;m\,(R_0+2r_0x_1)}{B^2L^2} \] % CHECK: 分子括号内写作 R_0，前面各式为 R
开始时和最后时 $I_{\min}=0$

设走过 $x$\quad $I=\dfrac{BL\sqrt{2ax}}{R+2r_0x}=\dfrac{BL\sqrt{2a}}{\dfrac{R}{\sqrt{x}}+2r_0\sqrt{x}}$\quad 均值不等式\quad $f(x)=\dfrac{1}{x+\dfrac{1}{x}}$

\[ \left\{\begin{array}{l}
\dfrac{R}{2r_0}<x_1\text{ 时}\quad I_{\max}=\dfrac{BL\sqrt{2a}}{2\sqrt{2Rr_0}}\\[3ex]
\dfrac{R}{2r_0}>x_1\text{ 时}\quad I_{\max}=\dfrac{BL\sqrt{2ax_1}}{R+2r_0x_1}
\end{array}\right. \]
%FIG p057_b 0.67 0.367 0.90 0.4755
\notefig[0.3]{figures/p057_b.png}
%FIG p057_c 0.575 0.475 0.78 0.55
\notefig[0.28]{figures/p057_c.png}
% 图 p057_c 中标注函数 x/(x^2+1)（奇函数曲线）
$\dfrac{1}{\sqrt{x}+\dfrac{1}{\sqrt{x}}}$
\end{solution}
\end{problem}
%%%END
%%%BLOCK id=057-02 ch=12 sec=电磁感应·动力学基本公式 kind=formula alt=none cont=none y=0.40-0.51
由法拉第电磁感应定律
\[ \left\{\begin{array}{l} E=BLV\\ I=\dfrac{E}{R_{\text{总}}}\\ F=BIL\end{array}\right.\qquad F=\frac{B^2L^2V}{R_{\text{总}}} \]
%%%END
%%%BLOCK id=057-03 ch=12 sec=杆类问题·分析要点 kind=method alt=none cont=none y=0.54-0.71
\textbf{最终速度、匀速、恒定}

① 最大速度 $F_{\text{合}}=0$\quad ② 加速度 $F_{\text{合}}=ma$\quad ③ 电荷量 $q=\dfrac{n\Delta\Phi}{R_{\text{总}}}$\qquad $\Delta q=I\Delta t=\dfrac{\Delta\Phi}{R}=\dfrac{BLx}{R}$

④ 功率 $P_{\text{电}}=P_{\text{安}}=F_{\text{安}}v=\dfrac{B^2L^2v^2}{R}$\quad ⑤ 热量：动能损失、$W_{\text{安}}$、$I^2Rt$（恒定 $I$）

⑥ 总时间总位移，$x_{\text{相}}$，$S$ 相对面积\quad 动量定理，动量守恒 $\sum F_{\text{外}}\cdot\Delta t-F_{\text{安}}\Delta t=m(V_{\text{末}}-V_0)-\Delta P$\quad $\left\{\begin{array}{l}BLq\\[0.5ex] \dfrac{B^2L^2x}{R}\end{array}\right.$

⑦ $L$ 不变，匀速，$B'\Delta S=B_0S_0$
\[ B'(t)=\frac{B_0S_0}{L\left(V_0t+\frac{1}{2}at^2-L_0\right)} \] % 此行左侧原稿有涂改液痕迹，残留一个逗号
%%%END
%%%BLOCK id=057-04 ch=12 sec=杆类问题·运动功能动量三方面 kind=summary alt=none cont=none y=0.71-0.80
% 原稿左侧竖写“\unclear{运动}”，用大括号括住 V_m 与 a；“P电”与“P安”竖排，中间为竖的等号
\[
\left.\begin{array}{l}V_m\\ a\end{array}\right\}\ \text{\unclear{运动}}
\]
功能\quad $\begin{array}{c}P_{\text{电}}\\ \|\\ P_{\text{安}}\end{array}$\quad $\left\{\begin{array}{l}W_{\text{除安}}-Q=\Delta E_k\\ P_{\text{电}}=P_{\text{安}}\\ W_{\text{电}}=W_{\text{安}}=-Q\end{array}\right.$ % CHECK: “W除安”的第二字辨认为“除”(W_除安 + W_安 = ΔE_k 即得此式)，第三行中间项为 W_安(原稿字迹偏潦草)

动量\quad $\left\{\begin{array}{l}q=n\dfrac{\Delta\Phi}{R}=n\cdot\dfrac{\Phi_2-\Phi_1}{R}\\[2ex] Ft-\dfrac{B^2L^2x}{R}=m(V_{\text{末}}-V_0)\end{array}\right.$\qquad 三方面
%%%END
%%%BLOCK id=057-05 ch=12 sec=电磁感应·并联电动势等效 kind=method alt=none cont=none y=0.80-0.93
%FIG p057_d 0.03 0.81 0.195 0.875
\notefig[0.25]{figures/p057_d.png}
2 个 $BdV$ 并联了\quad $e=BdV$\quad $I=\dfrac{BdV}{R}$\quad $F=\dfrac{B^2d^2V}{R}$
%FIG p057_e 0.20 0.845 0.345 0.93
\notefig[0.25]{figures/p057_e.png}
\[ e_{\text{等}}=BdV,\qquad r_{\text{等}}=\frac{r}{4} \]
%%%END
%%%BLOCK id=058-01 ch=12 sec=单杆·电势差与电路 kind=problem alt=none cont=none y=0.00-0.13
\begin{problem}[①]
%FIG p058_a 0.10 0.0 0.395 0.13
\notefig[0.4]{figures/p058_a.png}
单位长电阻 $r$，求 $U_{PQ}$
\begin{solution}
\[ U_{PQ}=2U_{\text{外}}+U_{mn}=\frac{Bdv}{R+rd}R+B(L-d)v=\frac{RL+rd\,(L-\text{\unclear{…}}}{R+rd} \] % 原稿末式分子在页面右缘被截断；CHECK: “2U_外”与后式系数是否一致待核
\end{solution}
\end{problem}
%%%END
%%%BLOCK id=058-02 ch=12 sec=线框·终极速度 kind=problem alt=none cont=none y=0.13-0.28
\begin{problem}[②]
%FIG p058_b 0.04 0.135 0.165 0.21
\notefig[0.25]{figures/p058_b.png}
\begin{solution}
% 原稿 V_m 结果 Dρg/B^2 用椭圆圈住
\[ \left\{\begin{array}{l} mg=\dfrac{B^2L^2V_m}{R}\\[2ex] m=DS\cdot 2\pi r\\[1ex] R=\rho\dfrac{2\pi r}{S}\end{array}\right.\qquad L=2\pi r\qquad \therefore V_m=\boxed{\dfrac{D\rho g}{B^2}} \]
%FIG p058_c 0.66 0.205 0.77 0.275
\notefig[0.25]{figures/p058_c.png}
\[ a=g-\frac{BV_0^2}{\text{\unclear{$16R\rho_2$}}} \] % CHECK: 分子指数 2 写在 V_0 之后，疑为 B^2V_0；分母 16Rρ_2 辨认不清(由 m=4LdA、R=ρ4L/A 推得应为 B^2V_0/(16dρ))
\end{solution}
\end{problem}
%%%END
%%%BLOCK id=058-03 ch=12 sec=单杆·恒力与恒功率启动 kind=problem alt=none cont=none y=0.27-0.48
\begin{problem}[③]
%FIG p058_d 0.07 0.275 0.195 0.36
\notefig[0.25]{figures/p058_d.png}
用 $F$ 拉 $ab$，$t$ 后 $ab$ 速度为 $V_1$，加速度 $a_1$，最终速度 $2V$。若拉力功率恒定，经 $t$ 后 $ab$ 速度 $V$ 加速度 $a_2$，最终速度 $2V$。求 $\dfrac{a_1}{a_2}$。 % V 右下角的小标记辨认为 1(亦可能是 V_A)
\begin{solution}
\[ \left\{\begin{array}{l}
F=\dfrac{B^2L^2\cdot 2V}{R}\\[2.5ex]
F-\dfrac{B^2L^2V}{R}=ma_1\\[2.5ex]
\dfrac{P}{2V}=\dfrac{B^2L^2\cdot 2V}{R}\\[2.5ex]
\dfrac{P}{V}-\dfrac{B^2L^2V}{R}=ma_2
\end{array}\right.\qquad \frac{a_1}{a_2}=\frac{1}{3} \]
\end{solution}
\end{problem}
%%%END
%%%BLOCK id=058-04 ch=12 sec=单杆·斜面导轨 kind=problem alt=06 cont=none y=0.48-0.66
\begin{problem}[④]
%FIG p058_e 0.04 0.515 0.24 0.585
\notefig[0.28]{figures/p058_e.png}
\unclear{释放}，到 $V$ 后匀速，加一沿斜面向下力，功率为 $P$，最终以 $2V$ 匀速（\unclear{棒}无电阻）\quad 求 $\dfrac{V}{2}$ 时的加速度、$P$、$2V$ 后的焦耳热
\begin{solution}
\[ \left\{\begin{array}{l}
mg\sin\theta=\dfrac{B^2L^2V}{R}\\[2.5ex]
\dfrac{P}{2V}+mg\sin\theta=\dfrac{B^2L^2\cdot 2V}{R}\\[2.5ex]
mg\sin\theta-\dfrac{B^2L^2\cdot\frac{V}{2}}{R}=ma\\[2.5ex]
W_G+W_F-W_{\text{安}}=0
\end{array}\right.\qquad \begin{array}{l}a=\dfrac{1}{2}g\sin\theta\\[2ex] P=2mg\sin\theta\,V\end{array} \] % 第四式首项原稿字形近似 W_a，按语境读作 W_G
\end{solution}
\end{problem}
%%%END
%%%BLOCK id=058-05 ch=12 sec=线框·下落终极速度与产热 kind=problem alt=06 cont=none y=0.66-0.86
\begin{problem}[⑤]
%FIG p058_f 0.055 0.69 0.16 0.76
\notefig[0.25]{figures/p058_f.png}
框 % 图旁标注
\begin{solution}
\[ \left\{\begin{array}{l}
mg=\dfrac{B^2(2L)^2V_m}{R}\\[2.5ex]
m=4L\cdot dA\\[1ex]
R=\rho\dfrac{4L}{A}
\end{array}\right. \]
\struck{$P_{\text{电}}=P_{\text{安}}V$}\qquad 当 $a=\dfrac{1}{2}g$ 时\quad $mg-\dfrac{B^2L^2V}{R}=m\dfrac{g}{2}$

$P_{\text{电}}=F_{\text{安}}V$。 % 原稿有一条大弧线箭头，自“当 a=g/2 时…”一行右上方指向被划掉的“P电=P安V”
\end{solution}
下落 $t$ 后、$h$、速为 $V_t$ $(<V_m)$ 在同 $t$ 内，框产热与恒定电流 $I_0$ 产热相同（在该框中），求 $I_0$。
\begin{solution}
\[ mgh-I_0^2Rt=\frac{1}{2}mV_t^2 \]
\end{solution}
\end{problem}
%%%END
%%%BLOCK id=058-06 ch=12 sec=单杆·匀加速运动分析 kind=method alt=none cont=none y=0.86-1.00
\textbf{匀加速}
\[ F-\frac{B^2L^2at}{R}=ma \]
%FIG p058_g 0.245 0.90 0.37 0.968
\notefig[0.25]{figures/p058_g.png}
\[ F=ma+\frac{B^2L^2a}{R}t \]
\[ q=\frac{BL\cdot\frac{1}{2}at^2}{R}\qquad I_{\text{安}}=BLq\qquad I=\frac{BLa}{R}t \] % q 的分母 R 被页面下缘截去一部分
%FIG p058_h 0.42 0.918 0.553 1.0
\notefig[0.25]{figures/p058_h.png}
\[ q=\frac{1}{2}I_1t_1\qquad\qquad W_F=Q_{\text{电}}+\frac{1}{2}m(at)^2 \]
%%%END
%%%BLOCK id=059-01 ch=12 sec=涡旋电场 kind=knowledge alt=none cont=none y=0.03-0.28
\textbf{涡旋电场}\quad \fbox{\unclear{非静电力}}\ 用有效磁场面积 % 原稿“非静电力”(三个字，字迹难辨)被椭圆圈起
%FIG p059_a 0.09 0.11 0.275 0.225
\notefig[0.3]{figures/p059_a.png}
$I$ $E_{\text{涡}}$ 同向\quad 方向楞次 Th
\[ E_{\text{涡}}=\left\{\begin{array}{ll}\dfrac{kr}{2}, & r<r_0\\[2ex] \dfrac{kr_0^2}{2r}, & r>r_0\end{array}\right. \]
%FIG p059_b 0.572 0.105 0.735 0.215
\notefig[0.28]{figures/p059_b.png}
\[ E=\frac{\dfrac{\Delta B}{\Delta t}\,\pi r^2}{2\pi r}=\frac{\dfrac{\Delta B}{\Delta t}\,r}{2}=\frac{kr}{2} \]
\[ E=\frac{\dfrac{\Delta B}{\Delta t}\,\pi r_0^2}{2\pi r}=\frac{kr_0^2}{2r} \]
\[ E=\frac{U}{d}\qquad U=\frac{\Delta B}{\Delta t}S\qquad d=2\pi r\ (\text{区域周长。})\qquad \underline{S\text{ 为包住磁场的面积}} \]
%%%END
%%%BLOCK id=059-02 ch=12 sec=旋转感应电动势 kind=knowledge alt=none cont=none y=0.27-0.38
\textbf{旋转感应电动势}
\[ e=BL\,\frac{V_1+V_2}{2} \]
$\uparrow$ \unclear{中点}瞬时速度 % 原稿箭头自 (V_1+V_2)/2 指向此文字，前两字重叠难辨
%FIG p059_c 0.325 0.283 0.435 0.35
\notefig[0.25]{figures/p059_c.png}
%FIG p059_d 0.462 0.27 0.655 0.375
\notefig[0.3]{figures/p059_d.png}
\[ S=\frac{\varepsilon}{B}=\frac{\omega}{2}(r_1+r_2)(r_2-r_1) \] % 页左边缘此处有一红笔括弧标记
%%%END
%%%BLOCK id=059-03 ch=12 sec=电容器与导体棒·恒力无电阻 kind=method alt=09 cont=none y=0.38-0.57
\textbf{电容与棒}\quad 电阻作用是缓冲

\textbf{恒力无电阻} % 原稿用椭圆圈起
%FIG p059_e 0.095 0.43 0.205 0.495
\notefig[0.25]{figures/p059_e.png}
\[ \left\{\begin{array}{l}
i=\dfrac{\Delta q}{\Delta t}=C\dfrac{\Delta U}{\Delta t}=C\dfrac{\mathrm{d}(BLV)}{\mathrm{d}t}=CBLa\\[2.5ex]
F_{\text{安}}=BiL\\[1ex]
a=\dfrac{F-BiL}{m}=\dfrac{F-CB^2L^2a}{m}
\end{array}\right. \]
\[ \therefore a=\frac{F}{CB^2L^2+m} \]
%FIG p059_f 0.465 0.515 0.58 0.58
\notefig[0.25]{figures/p059_f.png}
%%%END
%%%BLOCK id=059-04 ch=12 sec=电容器与导体棒·恒力有电阻 kind=method alt=09 cont=none y=0.40-0.75
\textbf{恒力有电阻} % 原稿用椭圆圈起
%FIG p059_g 0.605 0.46 0.715 0.54
\notefig[0.25]{figures/p059_g.png}
\[ i=C\frac{\Delta U}{\Delta t}=\frac{C\Delta(\varepsilon-U_R)}{\Delta t} \]
\[ =\frac{CBL\Delta V}{\Delta t}-\frac{C\Delta U_R}{\Delta t} \]
\[ =CBLa-C\frac{\Delta U_R}{\Delta t} \]
\[ F-B\left(CBLa-\frac{C\Delta U_R}{\Delta t}\right)L=ma \]
↓ 最终减为 0 % 原稿箭头自括号内 CΔU_R/Δt 指向“最终减为0”

稳定后 $\Delta\varepsilon=U_R+U_C$
%FIG p059_h 0.755 0.590 0.915 0.66
\notefig[0.28]{figures/p059_h.png}
\[ a=\frac{F}{CB^2L^2+m} \]
$\uparrow$ \unclear{视}在质量
%%%END
%%%BLOCK id=059-05 ch=12 sec=电容器与导体棒·初速加电阻 kind=method alt=09 cont=none y=0.57-0.76
\textbf{初速加电阻} % 原稿用椭圆圈起
%FIG p059_i 0.105 0.62 0.275 0.695
\notefig[0.25]{figures/p059_i.png}
\[ \left\{\begin{array}{l}
i=\dfrac{BLV-U_c}{R}\ \ \text{最终 }BLV=U_c\\[2.5ex]
F_{\text{安}}=BiL\quad a\downarrow\\[1ex]
mV_m-mV_0=-BLq\\[1ex]
U_c=BLV_m\\[1ex]
q=CU_c=CBLV_m
\end{array}\right.\qquad V_m=\frac{mV_0}{CB^2L^2+m} \]
%FIG p059_j 0.655 0.63 0.80 0.715
\notefig[0.25]{figures/p059_j.png}
%%%END
%%%BLOCK id=059-06 ch=12 sec=电容器与导体棒·初速无电阻 kind=method alt=09 cont=none y=0.76-0.85
\textbf{初速无电阻} % 原稿用椭圆圈起
%FIG p059_k 0.225 0.78 0.31 0.84
\notefig[0.25]{figures/p059_k.png}
%FIG p059_l 0.32 0.775 0.45 0.84
\notefig[0.25]{figures/p059_l.png}
\[ BiL=CB^2L^2a=ma \]
故 $a=0$
%%%END
%%%BLOCK id=059-07 ch=12 sec=自感·电流图像 kind=knowledge alt=none cont=none y=0.86-1.00
自感 $E=N\dfrac{\Delta I}{\Delta t}L$ % CHECK: 分子辨认为 ΔI(亦可能是 ΔΦ)，N 与末尾 L 的含义存疑，按原样转录
%FIG p059_m 0.33 0.868 0.475 0.93
\notefig[0.25]{figures/p059_m.png}
%FIG p059_n 0.295 0.928 0.45 0.995
\notefig[0.25]{figures/p059_n.png}
考虑自感
%FIG p059_o 0.55 0.895 0.745 0.99
\notefig[0.28]{figures/p059_o.png}
增减都阻碍
%%%END
%%%BLOCK id=060-01 ch=12 sec=线框·落框 kind=derivation alt=none cont=none y=0.04-0.24
\textbf{落框}
%FIG p060_a 0.125 0.045 0.28 0.115
\notefig[0.25]{figures/p060_a.png}
\[ V=\sqrt{2gh}\quad E=nBLV\quad m=D\,4nLS\quad R=\rho\frac{4nL}{S}\quad I=\frac{E}{R},\quad F_{\text{安}}=nBIL=\frac{mB^2V}{16D\rho} \]
\[ mg-F_{\text{安}}=ma\qquad a=g-\frac{F}{m}=g-\frac{B^2V}{16\rho D}\qquad \left\{\begin{array}{l} g>\dfrac{B^2V}{16D\rho}\ \text{时都加速}\\[2ex] g<\dfrac{B^2V}{16D\rho}\ \text{都减速}\end{array}\right. \] % CHECK: a=g-□/m 的分子字形介于 E 与 F，按物理意义读作 F(应为 F_安)
%%%END
%%%BLOCK id=060-02 ch=12 sec=线框·落框 kind=derivation alt=none cont=none y=0.20-0.38
%FIG p060_b 0.05 0.205 0.14 0.27
\notefig[0.2]{figures/p060_b.png}
\[ mg=\frac{B^2L^2V_y}{R}=\frac{B^2L^2\sqrt{2gH}}{R} \]
%FIG p060_c 0.025 0.29 0.21 0.366
\notefig[0.25]{figures/p060_c.png}
\[ \therefore B=\sqrt{\frac{mgR}{L^2\sqrt{2g}}}\ \frac{1}{\sqrt{\sqrt{H}}}\ \propto\ H^{-\frac{1}{4}} \]
\[ P_G=P_{\text{安}}\quad mg=F_{\text{安}}\quad W_{\text{热}}=F_{\text{安}}\times 2L\quad \text{进出匀速后，}F_{\text{安}}\text{不变} \]
%%%END
%%%BLOCK id=060-03 ch=12 sec=线框·穿越磁场的动量定理 kind=derivation alt=07 cont=none y=0.38-0.50
%FIG p060_d 0.095 0.41 0.24 0.49
\notefig[0.25]{figures/p060_d.png}
\[ -BLq=mV_2-mV_1\qquad \therefore\ |V_2-V_1|=|V_3-V_2| \]
\[ \Delta X=\frac{m\Delta VR}{B^2L^2}\qquad \left\{\begin{array}{l} q_1=q_2\\ Q_1=3Q_2\end{array}\right. \]
%%%END
%%%BLOCK id=060-04 ch=12 sec=线框·进出磁场v-t图分类讨论 kind=method alt=none cont=none y=0.49-0.72
\textbf{直接分类讨论}（与\unclear{束缚}洛伦兹力一样 直接分类讨论） % “束缚”二字辨认不清
%FIG p060_e 0.09 0.52 0.225 0.595
\notefig[0.25]{figures/p060_e.png}
%FIG p060_f 0.275 0.504 0.54 0.62
\notefig[0.3]{figures/p060_f.png}
% 图 p060_f 中标注：mg>F_安1、mg=F_安1、mg<F_安1 三种分支，以及起始段 mg>F_安
%FIG p060_g 0.55 0.51 0.75 0.62
\notefig[0.3]{figures/p060_g.png}
% 图 p060_g 中标注：入、出_1，并有批注“全等故产热相当”
不全等产热就不相同
%FIG p060_h 0.28 0.628 0.50 0.725
\notefig[0.25]{figures/p060_h.png}
% 图 p060_h 中标注：mg=F_安、进、出、mg<F_安1
%%%END
%%%BLOCK id=060-05 ch=12 sec=线框·落框进出磁场能量关系 kind=derivation alt=06 cont=none y=0.71-0.87
%FIG p060_i 0.08 0.715 0.26 0.84
\notefig[0.28]{figures/p060_i.png}
\[ V_{\text{入}}=V_{\text{出}} \]
\[ \left\{\begin{array}{ll}\text{进} & mg(h+L)-Q=\dfrac{1}{2}mV^2\\[2ex] \text{出} & mgd-Q=\dfrac{1}{2}mV^2-\dfrac{1}{2}mV^2\end{array}\right.\qquad \therefore mgd=Q \]
\[ V=\sqrt{2(h+L-d)g} \]
%%%END
%%%BLOCK id=060-06 ch=12 sec=线框·抛框上升与下降对比 kind=summary alt=none cont=none y=0.88-1.00
%FIG p060_j 0.13 0.893 0.292 0.97
\notefig[0.25]{figures/p060_j.png}
向上时的大速小时\ 大安力

等电荷量\ 等安冲
\[ \begin{array}{ll}
\overline{V}_{\text{上}}>\overline{V}_{\text{下}} & P_{\text{电上}}>P_{\text{电下}}\\
\overline{a}_{\text{上}}>\overline{a}_{\text{下}} & \varepsilon_{\text{上}}>\varepsilon_{\text{下}}\\
t_{\text{上}}<t_{\text{下}} & q_{\text{上}}=q_{\text{下}}\\
Q_{\text{上}}>Q_{\text{下}} & I_{\text{安冲上}}=I_{\text{安冲下}}
\end{array} \] % 原稿 t_下 上方似有一横线，疑为笔迹，未转录
%%%END
%%%BLOCK id=061-01 ch=12 sec=线框穿磁场·图像 kind=summary alt=none cont=none y=0.02-0.14
\textbf{线框穿磁场}\quad 考察右手定则、电源串联、有效长度 % CHECK: 原稿“穿”字小写在“传”字上方（疑为更正），主行写作“线框传磁场”
%FIG p061_a 0.39 0.03 0.57 0.068
\notefig[0.3]{figures/p061_a.png}
% 图中文字：相加、相减（两组电池串联符号）
\textbf{线框图像}\quad 同形状
%%%END
%%%BLOCK id=061-02 ch=12 sec=线框穿磁场·图像 kind=method alt=none cont=none y=0.14-0.31
\textbf{单场匀速}\quad 进出反向轴对称
%FIG p061_b 0.10 0.173 0.25 0.30
\notefig[0.3]{figures/p061_b.png}
% 图中文字：θ、v、sin x 形、L
\[ E=B\cdot v\tan\theta\,t\cdot v \]
%FIG p061_c 0.458 0.17 0.60 0.245
\notefig[0.3]{figures/p061_c.png}
%FIG p061_d 0.265 0.243 0.41 0.306
\notefig[0.3]{figures/p061_d.png}
% 图中文字：i、t、2L/v_0
\[ Q=\left(\frac{\dfrac{BLv}{\sqrt{2}}}{R}\right)^{2}\frac{2L}{v_0} \] % CHECK: 按原稿转录，式中未见电阻 R 因子（Q=I^2Rt 应有 R）
%%%END
%%%BLOCK id=061-03 ch=12 sec=线框穿磁场·图像 kind=method alt=none cont=none y=0.31-0.49
\textbf{反向场，两同夹异\ 中加倍}
%FIG p061_e 0.10 0.32 0.56 0.495
\notefig[0.55]{figures/p061_e.png}
% 图中文字：匀速、i、t（两个 i-t 图像）
%%%END
%%%BLOCK id=061-04 ch=12 sec=线框穿磁场·图像 kind=method alt=none cont=none y=0.51-0.62
\textbf{单一磁场匀加速}
%FIG p061_f 0.20 0.538 0.32 0.604
\notefig[0.25]{figures/p061_f.png}
%FIG p061_j 0.335 0.527 0.49 0.611
\notefig[0.3]{figures/p061_j.png}
% 图中文字：a、i、t
前后对称后共线
%%%END
%%%BLOCK id=061-05 ch=12 sec=线框穿磁场·图像 kind=method alt=none cont=none y=0.61-0.74
\textbf{反向场匀加速}
%FIG p061_g 0.24 0.632 0.335 0.70
\notefig[0.3]{figures/p061_g.png}
%FIG p061_h 0.34 0.612 0.482 0.74
\notefig[0.3]{figures/p061_h.png}
两同夹异中加倍，（前后共线\ 中加倍）\\
前后图象共有线
%%%END
%%%BLOCK id=061-06 ch=12 sec=线框穿磁场·图像 kind=method alt=none cont=none y=0.74-0.87
%FIG p061_i 0.13 0.74 0.60 0.86
\notefig[0.6]{figures/p061_i.png}
% 图中文字：v→、L、2、i、t、进、过 L/2、未到界线之已进、往外出之往里进（图中有一处被修正带涂白）
图中批注：「进」「过$\dfrac{L}{2}$」；\unclear{未到界线}之已进；往外出之往里进 % CHECK: “未到界线”“往外”处字迹难辨
%%%END
%%%BLOCK id=083-15 ch=12 sec=电磁感应·图象 kind=method alt=01 cont=none y=0.71-0.81
% 原稿“□→”为一个小方块加向右箭头
$\square\to$\quad $\varepsilon=Bl\sqrt{2ax}$\quad \red{匀加速}
%FIG p083_f 0.335 0.71 0.52 0.81
\notefig[0.3]{figures/p083_f.png}
%%%END
%%%BLOCK id=097-04 ch=12 sec=磁通量 kind=formula alt=none cont=none y=0.22-0.26
$\Delta\phi=\phi_{\text{末}}-\phi_{\text{初}}$\qquad 面要分正负\qquad $\phi=BS\cos\theta$

%FIG p097_c 0.67 0.205 0.79 0.262
\notefig[0.25]{figures/p097_c.png}
%%%END
%%%BLOCK id=097-05 ch=12 sec=感应电荷量 kind=formula alt=none cont=none y=0.27-0.37
%FIG p097_d 0.11 0.265 0.33 0.375
\notefig[0.3]{figures/p097_d.png}

转过$2\pi$\ $q=\dfrac{4BS}{R}$，转过\unclear{$\dfrac{7}{4}\pi$}\ $q=\left(3+\dfrac{\sqrt{2}}{2}\right)\dfrac{BS}{R}$

$q$要按过程求

$q=n\dfrac{\Delta\phi}{\Delta t}$ % CHECK: 手写分母似为 Δt（通常 q=nΔΦ/R）；另“转过 7π/4”的分数辨认不确定
%%%END
%%%BLOCK id=097-06 ch=12 sec=楞次定律 kind=knowledge alt=none cont=none y=0.40-0.62
% 原稿“楞次定律”标题被菱形框框起
\textbf{楞次定律}

回路中\uwave{感应电流总阻碍}\uwave{原磁通量变化}（不是阻止，来自于能量守恒）

% 原稿以一个大括号把下面四行括在一起
\begin{itemize}
\item $\pm$增反减同
\item 来拒去留
\item 线圈靠磁铁\quad $\pm$增扩减缩
\item 线圈不靠磁铁\quad $\pm$增缩减扩
\end{itemize}

二次感应，一样，$i_{\text{感}}=\dfrac{\dd i_{0}}{\dd t}$\quad 导数关系。 % CHECK: 分子下标手写似“0/o”，可能是“原”（原电流）
%%%END
%%%BLOCK id=097-07 ch=12 sec=感生电场·图像选择题 kind=method alt=none cont=none y=0.65-0.80
\textbf{感生电场选择题}（没有正方向，则自己规定）

$B_{\text{正}}\ I_{\text{正}}$\quad 满足右手，斜正电负，斜负电正（否则相反）

$F_{\text{正}}\ B_{\text{正}}\ I_{\text{正}}$\quad 满足右手，斜正对称，斜负相同（否则相反）

$B$、$F$关于$X$轴对称
% 原稿中“电负”与“对称”被圈在一起，“电正”与“相同”被圈在一起，第一行的“（否则相反）”也被圈起
%FIG p097_e 0.635 0.645 0.78 0.81
\notefig[0.3]{figures/p097_e.png}
%%%END
%%%BLOCK id=097-08 ch=12 sec=感生电动势·线框受力 kind=derivation alt=none cont=none y=0.82-0.93
%FIG p097_f 0.06 0.82 0.16 0.925
\notefig[0.25]{figures/p097_f.png}

\[ i=\dfrac{\dfrac{\Delta B}{\Delta t}\cdot\dfrac{L^{2}}{2}}{\dfrac{\rho\,4L}{S}}\qquad
\dfrac{\Delta F}{\Delta t}=\dfrac{\Delta B}{\Delta t}IL=KIL \]
\[ F=KILt \]
%%%END
%%%BLOCK id=098-04 ch=12 sec=电磁感应·动量定理 kind=problem alt=07,03 cont=none y=0.44-0.56
③
%FIG p098_c 0.04 0.447 0.215 0.56
\notefig[0.25]{figures/p098_c.png}

经$t_0$停，求$x$。由动量定理
\[ mg\sin\theta\,t_0+\frac{B^2\ell^2x}{R}=mV \]
% CHECK: 右端手写为 mV（无下标），按后式推算似应为 mV_0
\[ x=\frac{(V_0-gt_0\sin\theta)\,mR}{B^2\ell^2} \]
%%%END
%%%BLOCK id=099-05 ch=12 sec=电磁感应·变长度棒与1/v-x图像 kind=problem alt=01,06 cont=none y=0.55-0.80
⑤
%FIG p099_f 0.025 0.572 0.23 0.668
\notefig[0.28]{figures/p099_f.png}

%FIG p099_g 0.24 0.563 0.445 0.68
\notefig[0.28]{figures/p099_g.png}

由$\dfrac{1}{V}\propto x$\ 为\uline{非匀变速运动}

\red{\uwave{图像面积为运动时间}}\quad $t=\dfrac{3L_0}{2V_0}$
\[ \frac{1}{V}=\frac{1}{L_0V_0}\cdot x+\frac{1}{V_0}\ \to\ V=\frac{L_0V_0}{x+L_0} \]
\[ E=B(L_0+x)V\qquad I=\frac{E}{R}=\frac{BL_0V_0}{R}\ \text{（定值）} \]
\[ W=I^2Rt-\Delta E_k=\frac{3B^2L_0^2V_0}{2R}-\frac{3mV_0^2}{8} \]
% CHECK: 分子中 L_0 的指数手写为 2，按 I^2Rt 推算似应为 L_0^3
%%%END
%%%BLOCK id=115-01 ch=12 sec=动生与感生电动势叠加 kind=method alt=none cont=none y=0.04-0.16
动生$+$感生\ ·\ (对$t$求导)\unclear{相}

$\varepsilon=\dfrac{\Delta B}{\Delta t}\cdot S+BLv$

\begin{itemize}
\item 给 $B=kt$ 直接加\quad $e=BLv+kS$\quad $F_{\text{安}}=BL\,\dfrac{BLv+kS}{R}$
\item 没给 则按电流方向是否相同来加减
  \begin{itemize}
  \item 同向$\to$相加
  \item 反向$\to$相减．
  \end{itemize}
\end{itemize}
%%%END
%%%BLOCK id=115-02 ch=12 sec=动生与感生电动势叠加 kind=problem alt=none cont=none y=0.16-0.50
% 原稿此处贴有印刷体题目（题号B4），其中若干短语被手画方框，“杆以恒定”下有横线；序号①为圈起的手写“1”
\begin{problem}[1（B4）]
B4、如图所示，两根平行金属导轨固定在水平桌面上，每根导轨每米的电阻为\fbox{$r_0=0.10\unit{\Omega/m}$,}导轨的端点 P、Q 用电阻可忽略的导线相连，两导轨间的距离\fbox{$L=0.20\unit{m}$.}有随时间变化的匀强磁场垂直于桌面，已知磁感应强度 $B$ 与时间 $t$ 的关系为\fbox{$B=kt$}\ 比例系数 $k=0.020\unit{T/s}$。一电阻不计的金属杆可在导轨上\fbox{无摩擦地滑动}，在滑动过程中保持与导轨垂直。在 $t=0$ 时刻，金属杆紧靠在 P、Q 端，在外力作用下，\underline{杆以恒定}的\fbox{加速度}\fbox{从静止开始}向导轨的另一端滑动，求\fbox{在 $t=6.0\unit{s}$ 时}金属杆所受的\fbox{安培力}。

% 图p115_a为印刷图（含P、Q、L、箭头、I），上有手写的 ½at² 与 x 标注
%FIG p115_a 0.095 0.383 0.40 0.452
\notefig[0.4]{figures/p115_a.png}
\begin{solution}
\[\left\{\begin{array}{l}S=Lx=\dfrac{1}{2}at^2\cdot L\\[1mm]B=kt\\[1mm]V=at\end{array}\right.\qquad I=\frac{BLat+kL\cdot\frac{1}{2}at^2}{2\cdot\frac{1}{2}at^2r_0}=\frac{3kL}{2r}\]
\[F_{\text{安}}=kt\cdot\frac{3kL}{2r}\cdot L=\frac{3k^2L^2}{2r}\,t=1.44\times10^{-3}\unit{N}\qquad\text{\struck{$=\dfrac{3kt}{2r}$}}\]
\end{solution}
\end{problem}
%%%END
%%%BLOCK id=115-03 ch=12 sec=电磁感应中的能量·收尾速度 kind=problem alt=none cont=none y=0.47-0.85
\begin{problem}[2]
% 图p115_b内含标注：d、m、R（金属环与磁场示意）
%FIG p115_b 0.07 0.513 0.215 0.588
\notefig[0.25]{figures/p115_b.png}

金属环下落，求收尾速度\quad \unclear{目}重力势能$\to$电能

$B_y=B_0(1+ky)\ (k>0)$
\begin{solution}
由功能关系\qquad 合外力为0\ $\times$ % 原稿此处打叉
\[mgV_m=P_{\text{电}}=\frac{\varepsilon^2}{R}\]
\[\varepsilon=\frac{\dd\left(B_0(1+ky)\dfrac{\pi d^2}{4}\right)}{\dd t}=B_0\frac{\pi kd^2}{4}V_m\]
\red{$\uparrow$ 对上式对$t$求导}
\[\therefore\ V_m=\frac{16mgR}{B_0^2\pi^2k^2d^4}\]
\underline{求导过程}
\[\frac{\dd\left(B_0(1+ky)\dfrac{\pi d^2}{4}\right)}{\dd t}=\frac{B_0\,k\,v\,\pi d^2}{4}\]
\end{solution}
\end{problem}
%%%END
%%%BLOCK id=137-01 ch=12 sec=电磁感应中的动量 kind=derivation alt=none cont=none y=0.04-0.27
% 原稿：几个式子被手绘的大圈/小圈圈住，并有一条斜线从圈引向右上
\[ -B\bar{I}L\,\Delta t=mv_2-mv_1 \]
\[ -BLq=mv_2-mv_1 \]
\[ q=n\,\frac{\Delta\varphi}{\Delta t} \]% CHECK: 分母原稿看似 Δt（常见写法为 R）
\[ -\frac{B^2L^2\bar{v}}{R}\,\Delta t=\Delta p \]
\[ \Delta x=\frac{m\,\Delta v\,R}{B^2L^2} \]
%FIG p137_a 0.54 0.075 0.72 0.16
\notefig[0.25]{figures/p137_a.png}
\[ |v_2-v_1|=|v_3-v_2| \]
\[ q_{\text{进}}=q_{\text{出}} \]
%%%END
%%%BLOCK id=158-01 ch=12 sec=电磁感应·磁单极子穿线圈 kind=knowledge alt=none cont=none y=0.05-0.34
% 三行图：左为磁铁(S N)/磁单极子(圈N)穿过线圈的示意(标有“斥”“引”“i感max”等)，右为对应的 i-t 图(标“中点0”)；图内文字已在图中
\noindent 〈磁单极子〉穿框　$i$方向不变

\noindent 超导线圈　$R=0$　$I$不会随时间衰减.
%FIG p158_a 0.05 0.133 0.40 0.22
\notefig[0.35]{figures/p158_a.png}
%FIG p158_b 0.05 0.21 0.40 0.285
\notefig[0.35]{figures/p158_b.png}
%FIG p158_c 0.05 0.275 0.40 0.34
\notefig[0.35]{figures/p158_c.png}
%%%END
%%%BLOCK id=158-02 ch=12 sec=电磁感应·法拉第圆盘 kind=knowledge alt=none cont=none y=0.11-0.40
% “盘转”等处的“盘”字形与标题“法拉第圆盘”的“盘”相同；“涡流”“磁场”之间上方有一个小字，难辨
\noindent \textbf{法拉第圆盘}
%FIG p158_d 0.53 0.135 0.675 0.225
\notefig[0.25]{figures/p158_d.png}
\[ I=\dfrac{BL^2\omega}{2R} \]
\noindent 若$B$改为$B\sin\omega t$，不转盘则产生涡流，灯不会亮（涡流不能流出）。
%FIG p158_e 0.52 0.25 0.625 0.31
\notefig[0.25]{figures/p158_e.png}
\noindent 盘转磁针也转

\noindent 整个中不变　局部变化

\noindent 没有自由电子形成电流

\noindent 涡流\textsuperscript{\unclear{…}}磁场使磁针转动
%%%END
%%%BLOCK id=158-03 ch=12 sec=感生电场·涡旋电场 kind=knowledge alt=09,04 cont=none y=0.35-0.58
% 原稿：“做功W=qES”起至“转圈数”一行左侧有一个大括号形曲线括起；W_1圈 一行在括号之外；“电子加速器只能加速电子”末尾“电子”二字下有横线；“顺逆均可”一行上方有小字“顺逆时针”
\noindent 〈感生电场〉　耗散力　加速圆周运动　$a$不指向圆心

\noindent 做功　$W=qES$

\noindent \hspace*{2em}$S=V_0t+\dfrac{1}{2}\dfrac{qE}{m}t^2$

\noindent \hspace*{2em}$V=V_0+at$

\noindent 转圈数　$n=\dfrac{\frac{1}{2}mV_0^2}{e\,E_{\text{感应电动势}}}=\dfrac{\text{末动能}}{\text{做功/圈}}$% CHECK: 分子 ½mV_0^2 的下标0存疑(末动能应为 ½mV^2)

\noindent $W_{1\text{圈}}=Uq=\dfrac{\Delta B}{\Delta t}\pi r^2q=k\pi r^2q$

\noindent 电子加速器只能加速\underline{电子}

\noindent 顺逆均可，\textsuperscript{顺逆时针}做功不相等.
%%%END
%%%BLOCK id=159-01 ch=12 sec=电磁感应·能量与电量 kind=problem alt=none cont=prev y=0.00-0.38
\begin{problem}
%FIG p159_a 0.16 0.00 0.41 0.135
\notefig[0.3]{figures/p159_a.png}
图旁标注已知量：$B=0.4\unit{T}$，$L=0.5\unit{m}$，$m=0.1\unit{kg}$，$v_0=2\unit{m/s}$，匀减速 $a=2\unit{m/s^2}$，$\mu=0$，棒无电阻，框电阻不均匀。

① 导棒运动0.5s，框架产热 % CHECK: 此处“0.5s”也许指位移0.5m（由 v^2-v_0^2=-2as 得 Q=0.1J 才自洽；若为0.5s则 Q=0.15J）
\begin{solution}
\[ v^2-v_0^2=-2as \]
\[ Q=\tfrac12mv_0^2-\tfrac12mv^2=0.1\unit{J} \]
\end{solution}
② 求$R(x)$及0.4s内通过棒的电量
\begin{solution}
\[ v=\sqrt{v_0^2-2a(x_0-x)}\qquad R=0.4\sqrt{x} \]
\[ q=It \]
\[ F=BIL=ma \]
\[ \therefore q=\frac{ma}{BL}\,t=0.4\unit{C} \]
（旁注）不能把0.4s时的电阻代入$q=\dfrac{BLx}{R}$
\end{solution}
\end{problem}
%%%END
%%%BLOCK id=159-02 ch=12 sec=电磁感应·动量定理与多磁场区 kind=problem alt=03 cont=none y=0.37-0.90
% 本页未见这道题的完整题干（已知条件），只有图（图中标注“穿过n个磁场”）和两问
\begin{problem}
%FIG p159_b 0.105 0.385 0.43 0.53
\notefig[0.4]{figures/p159_b.png}
① 导棒能匀速过$\mathrm{I}$，求释放地距$\mathrm{I}$距离 % CHECK: “释放地”的“地”字疑为“点”
\begin{solution}
\[ I=\frac{E}{R+\frac{R}{2}}\qquad E=BLv\qquad F_{\text{安}}=\frac{2B^2L^2v}{3} \] % CHECK: 原稿 F安 的分母只写了3，未见 R（按 I=E/(3R/2) 应为 3R）
\[ mg\sin30^\circ=F_{\text{安}} \]
\[ mgx_1\sin30^\circ=\tfrac12mv^2 \]
\[ \therefore x_1=\frac{9gm^2R^2}{16B^4L^4} \]
\end{solution}
② 导棒在相邻磁场间的无磁场运动的时间都为$t_0$（$n=1,2,\cdots$），求ab释放距离$\mathrm{I}$$x_2$及过每个磁场时间$t$
\begin{solution}
设棒进场速度 $v_1,v_3,\cdots,v_{2n-1}$，出磁场 $v_2,v_4,\cdots,v_{2n}$
\[ mg\sin\theta=ma \]
\[
\left\{\begin{aligned}
&d=\frac{v_2+v_3}{2}t_0=\frac{v_4+v_5}{2}t_0\\
&v_3-v_2=v_5-v_4=at_0\\
&\therefore v_3=v_5=\cdots=v_{2n-1}\\
&\phantom{\therefore\ }v_2=v_4=\cdots=v_{2n}
\end{aligned}\right.
\]
\[
\left\{\begin{aligned}
&\mathrm{I}\to\mathrm{II}\quad d=v_2t_0+\tfrac12at_0^2\\
&v_1=v_3=v_2+at_0\\
&v_1^2=2ax_2
\end{aligned}\right.
\]
\[ x_2=\frac{(4d+gt_0^2)^2}{16gt_0^2} \]
由动量定理
\[ (mg\sin30^\circ-F_{\text{安}})t=mv_2-mv_1 \]
\[ F_{\text{安}}t=\frac{2B^2L^2S}{3R} \]
\[ \therefore t=\frac{4SB^2L^2}{3mgR}-t_0 \]
\end{solution}
\end{problem}
%%%END
%%%BLOCK id=159-03 ch=12 sec=电磁感应·双棒模型 kind=derivation alt=none cont=none y=0.80-0.93
\textbf{求导：}
\[ (e=BLv_1-BLv_2)' \]
\[ 0=BLa_1-BLa_2 \]
\[ \therefore a_1=a_2 \]
%%%END
%%%BLOCK id=160-01 ch=12 sec=电磁感应·双棒模型 kind=problem alt=06 cont=none y=0.00-0.67
% 题干是粘贴的印刷体；印刷体上有手写划线（"导轨光滑且电阻忽略不计""两根质量均为m、有效电阻均为R的导体棒a和b放在导轨上"一带有下划线）及几处连线/斜杠记号，已略
% 图上方有手写的 L（轨距）记号，在图的裁剪范围内
\begin{problem}
（题干首行被页缘截断，仅见后半）\unclear{足够}长的平行金属导轨与水平面的夹角为$\theta$，导轨光滑且电阻忽略不计。场强为$B$的条形匀强磁场方向与导轨平面垂直，磁场区域的宽度为$d_1$，间距为$d_2$。两根质量均为$m$、有效电阻均为$R$的导体棒a和b放在导轨上，并与导轨垂直，重力加速度为$g$。

(1) 若a进入第2个磁场区域时，b以与a同样的速度进入第1个磁场区域，求b穿过第1个磁场区域过程中增加的动能$\Delta E_k$；

(2) 若a进入第2个磁场区域时，b恰好离开第1个磁场区域；此后a离开第2个磁场区域时，b又恰好进入第2个磁场区域。且a、b在任意一个磁场区域或无磁场区域的运动时间均相等。求b穿过第2个磁场区域过程中，两导体棒产生的总焦耳热$Q$；

(3) 对于第(2)问所述的运动情况，求a穿出第$k$个磁场区域时的速率$v_k$。

%FIG p160_a 0.04 0.35 0.56 0.60
\notefig[0.5]{figures/p160_a.png}
\begin{solution}
(1)\quad $\therefore\ \Delta E_k=mgd_1\sin\theta$

(2)
\[ \tfrac12mv_1^2+Q=\tfrac12mv_2^2+mgd_1\sin\theta \]
\[ \tfrac12mv_2^2=\tfrac12mv_1^2+mgd_2\sin\theta \] % CHECK: 原稿此行末的 θ 写得像 α
\[ Q=mg(d_1+d_2)\sin\theta \]
（写在印刷题干(2)与(3)之间的批注）$I_{\text{安}}$\unclear{也}在每个磁场中不变，故$\Delta v_1=\Delta v_2=\Delta v_3$，形成速度交替

（以下一组算式写在图的下方，对应第(3)问，原稿无题号）
\[
\left\{\begin{aligned}
&m(v_1-v_2)=mgt\sin\theta-\frac{B^2L^2d_1}{2R}\\
&d_2=\frac{v_1+v_2}{2}t\\
&v_1-v_2=-gt\sin\theta
\end{aligned}\right.
\]
\[ v_1=\frac{4mgRd_2}{B^2L^2d_1}\sin\theta-\frac{B^2L^2d_1}{8mR} \]
\end{solution}
\end{problem}
%%%END
%%%BLOCK id=162-01 ch=12 sec=电磁感应·线框进出磁场 kind=knowledge alt=none cont=none y=0.05-0.27
% 首行标题模糊
\unclear{线圈过程中}\quad $n$匝\unclear{金线}\ \unclear{$N_2$}

$V_0$与$\Delta V$无关\qquad $a=\dfrac{B^2V_0}{16\rho_1\rho_2}$\quad $m=4\rho_1LS$\quad $R=\rho_2\dfrac{4L}{S}$\quad $q=\dfrac{BLS}{4\rho_2}=\dfrac{BL^2}{R}$

$V_0\propto\Delta E_k$\qquad \unclear{$a$、$\rho_1$与所…有关} % 此行右半极模糊，只辨出个别字

$m\propto\Delta E_k$

$m\propto\dfrac{1}{\Delta V}$

$\Delta V=\dfrac{-B^2L^2\Delta x}{mR}$

$\Delta E_k=W_{\text{安}}$
%%%END
%%%BLOCK id=162-02 ch=12 sec=电磁感应·安培力冲量 kind=note alt=none cont=none y=0.17-0.24
若匀变速\ $F=kt$
%FIG p162_b 0.69 0.175 0.81 0.235
\notefig[0.25]{figures/p162_b.png}
$I_F=S_{\text{面积}}$
%%%END
%%%BLOCK id=162-03 ch=12 sec=电磁感应·能量与电量 kind=problem alt=06 cont=none y=0.29-0.52
% 本页未见这道题的文字题干，只有图和已知量，及解答算式
\begin{problem}
%FIG p162_a 0.075 0.295 0.36 0.395
\notefig[0.4]{figures/p162_a.png}
$B=2\unit{T}$，$R=3\,\Omega$，$r=1\,\Omega$，$L=1\unit{m}$，$m=2\unit{kg}$，$\mu=0.5$
\begin{solution}
\[ v=2x\qquad F_A=\frac{B^2L^2\cdot2x}{R+r}=2x\ (\mathrm{N}) \]
%FIG p162_c 0.48 0.375 0.605 0.43
\notefig[0.2]{figures/p162_c.png}
\[ W_A=-\tfrac12\,2\times1=-1\unit{J} \]
\[ W-\mu mgx_1+W_A=\tfrac12mv_1^2\qquad W=15\unit{J} \]
\[ q_x=\frac{\Delta\Phi}{R+r}=\frac{BLx}{R+r}=0.5\unit{C} \]
\[ Q=-W_{\text{安}}=1\unit{J}\qquad Q_R=\frac{R}{R+r}Q=0.75\unit{J} \]
\end{solution}
\end{problem}
%%%END
%%%BLOCK id=162-05 ch=12 sec=电磁感应·图像 kind=method alt=none cont=none y=0.63-0.78
%FIG p162_e 0.05 0.64 0.20 0.77
\notefig[0.25]{figures/p162_e.png}
\[ a=-\frac{B^2L^2V}{mR} \]
\[ I=\frac{BLV}{R}\propto V \]
\[ \frac{\Delta V}{\Delta x}=-\frac{B^2L^2}{mR} \]
%FIG p162_g 0.42 0.635 0.76 0.78
\notefig[0.5]{figures/p162_g.png}
%%%END
%%%BLOCK id=162-06 ch=12 sec=电磁感应·线框穿过磁场区 kind=derivation alt=01 cont=none y=0.745-1.00
%FIG p162_f 0.04 0.752 0.375 0.87
\notefig[0.4]{figures/p162_f.png}
$t_1\sim t_2$匀速\quad$F_{\text{安}}=mg$\quad$\dfrac{B^2L^2V_1}{R}=mg$\quad$L=V_1(t_2-t_1)$\quad$B^2=\dfrac{mgR}{(t_2-t_1)V_1^3}$ % CHECK: 代入 L=V_1(t_2-t_1) 后分母应含 (t_2-t_1)^2；原稿看不出平方
\[ I_{\max}=\frac{BLV_2}{R} \]
\[ F_{\text{安}\max}=\frac{B^2(t_2-t_1)^2V_1V_2}{R}=\frac{mgV_2}{V_1}\qquad P_{\text{安}\max}=F_{\text{安}\max}V_2=\frac{mgV_2^2}{V_1} \] % CHECK: 分子 V_1V_2 处按推导应为 V_1^2V_2
%%%END
%%%BLOCK id=163-01 ch=12 sec=电磁感应·线框落入磁场比较 kind=method alt=03,06 cont=none y=0.07-0.33
%FIG p163_a 0.02 0.07 0.345 0.225
\notefig[0.4]{figures/p163_a.png}
\begin{tabular}{@{}l@{\qquad}l@{}}
①\,若②加速或匀速 & ①加速\\
\hphantom{①\,}若②减速 & ①加、减、匀均可
\end{tabular}

% 下图（花括号内三种情形）图内文字：②加速 a=g-B^2L^2v/m（v-t 曲线上升并趋近虚线）；②减速（v-t 曲线下降并趋近虚线）；③匀速（水平线）
% CHECK: a=g-B^2L^2v/m 的分母只见 m，似缺 R；第三项序号写得像③（疑为②匀速）
%FIG p163_b 0.78 0.11 0.985 0.265
\notefig[0.3]{figures/p163_b.png}
图象注意起点

①②在磁场中位移相同

$V_{2\text{入}}$小，$F_{\text{安}2}$小，$W_{F_{\text{安}2}}$小，$Q_2<Q_1$
\[ mgh-Q=\tfrac12mv^2\qquad \text{故}\ V_2>V_1 \]
\[ mgt-\frac{2B^2L^2L'}{R}=mV\qquad \text{由}\ V_2>V_1\ \Rightarrow\ t_1<t_2 \]
%%%END
%%%BLOCK id=163-02 ch=12 sec=电磁感应·方法提示 kind=note alt=none cont=none y=0.31-0.36
\unclear{函}数类直接代入

弹簧＋电磁：阻尼振动\quad 能量损耗
%%%END
%%%BLOCK id=163-04 ch=12 sec=电磁感应·弹簧线框 kind=problem alt=06,03 cont=none y=0.40-0.57
% 本页未见这道题的文字题干，只有图和解答算式
\begin{problem}
%FIG p163_d 0.02 0.40 0.20 0.535
\notefig[0.25]{figures/p163_d.png}
原长释放\quad $K=\dfrac{mg}{L}$
\begin{solution}
\[ mgt-I_T=mV \]
\[ 0.5L\neq\tfrac12Vt\ \text{非匀变速}\qquad I_T\neq\frac{mgL}{V}-mV \]
\[ E=BLV\qquad U_{ad}=-\frac{E}{R}\cdot\frac34R=-\frac34BLV \]
由牛二
\[ \left|mg-K\frac{L}{2}-BIL\right|=ma\qquad I=\frac{BLV}{R}\qquad a=\left|\frac{g}{2}-\frac{B^2L^2V}{mR}\right| \]
静止时 $Kx'=mg$

$\therefore x'=L$\qquad $E_p=\tfrac12Kx'^2=0.5mgL$

$\Delta E_p=mgL$

由能量守恒\quad $W_G-W_T=Q$\qquad $Q=0.5mgL$
\end{solution}
\end{problem}
%%%END
%%%BLOCK id=163-05 ch=12 sec=电磁感应·变有效长度 kind=problem alt=06 cont=none y=0.56-0.78
% 本页未见这道题的文字题干，只有图和解答算式
\begin{problem}
%FIG p163_e 0.009 0.565 0.31 0.675
\notefig[0.4]{figures/p163_e.png}
$P_R$不变\ 向左移$L$过程
\begin{solution}
位移$x$\quad $l=L-\dfrac{x}{2}$
\[ B\,l\,V=BLV_0\qquad \therefore V=\frac{2LV_0}{2L-x}\qquad x\uparrow\ V\uparrow\ \text{非匀变速} \] % 原稿“l”与“L”字形相近，按物理意义区分
\[ \Delta t=\frac{\Delta\Phi}{BLV_0}=\frac{B\frac{L}{2}\left(L+\frac{L}{2}\right)}{BLV_0}=\frac{3L}{4V_0} \]
\[ Q=I^2R\Delta t=\left(\frac{BLV_0}{R}\right)^2R\Delta t=\frac{3B^2L^3V_0}{4R} \]
$V=2V_0$时\quad $W=\tfrac12m(2V_0)^2-\tfrac12mV_0^2+Q=\dfrac{3B^2L^3V_0}{4R}+\dfrac32mV_0^2$

（写在“$V=2V_0$时”下方的小字）\unclear{即}（$x=L$时）
\end{solution}
\end{problem}
%%%END
%%%BLOCK id=164-01 ch=12 sec=电磁感应·双棒模型 kind=problem alt=03,06 cont=none y=0.00-0.43
% 本页未见这道题的文字题干，只有图和解答；顶部图被页边切掉了上部
% 原稿“l”与“L”字形相近，按所见照录
\begin{problem}
%FIG p164_a 0.17 0.00 0.485 0.09 soft
\notefig[0.4]{figures/p164_a.png}
先放1后放2
\begin{solution}
\[ E_1=BLv_0\qquad i=\frac{E_1}{2R}\qquad F_{\text{安}}=\frac{B^2l^2v_0}{2R} \]
① 若 $mg\sin\theta>\dfrac{B^2l^2v_0}{2R}$

\quad 对1\ $a_1=g\sin\theta-\dfrac{B^2l^2v_0}{2mR}$

\quad 对2\ $a_2=g\sin\theta+\dfrac{B^2l^2v_0}{2mR}$

%FIG p164_b 0.57 0.15 0.78 0.235
\notefig[0.25]{figures/p164_b.png}
$v_{\text{相}}\downarrow$\quad $E=BLv_{\text{相}}\downarrow$\quad $I\downarrow$\quad $F_{\text{安}}\downarrow$\quad $a_1=a_2$\quad $\phi$不变\quad $v_1=v_2$\quad $E_{\text{机}}$先$\downarrow$后不变

（写在“$F_{\text{安}}\downarrow$”下方）$F_{\text{安}}\to0$

② 若 $mg\sin\theta<\dfrac{B^2l^2v_0}{2R}$\quad $a_1$向上\quad $a$先$\downarrow$后$\uparrow$(向下)\quad 最终$v_1=v_2$\quad $\phi$定值\quad $W_{F_{\text{安}2}}>0$

%FIG p164_c 0.12 0.335 0.33 0.43
\notefig[0.25]{figures/p164_c.png}
$v_1>v_2$\quad $W_{F_{\text{安}1}}<0$ % 原稿此式写在“$v_1>v_2$”同行右侧

$W_{F_{\text{安}1}}>W_{F_{\text{安}2}}$\quad $E_{\text{机}}$先$\downarrow$后不变 % CHECK: $W_{F安1}<0$ 与 $W_{F安1}>W_{F安2}$ 并列，疑为比较绝对值
\end{solution}
\end{problem}
%%%END
%%%BLOCK id=196-05 ch=12 sec=电磁感应·恒功率拉框与动能定理 kind=problem alt=03,06 cont=none y=0.43-0.98
\begin{problem}
\notefig[0.4]{figures/p196_b.png}

\noindent $m=0.4\unit{kg}$\\
$R=3\,\Omega$\\
$L=1\unit{m}$\\
$\mu=0$\\
$B=1\unit{T}$

用恒功率为$27\unit{W}$的力拉框，在$1$处$MN$以$V_0$（未知）匀速通过磁场，当$MN$过$2$时，在导体框上再施加一水平向左的外力$F'$，外力大小为$F'=kV$\quad $k=\dfrac{4}{3}\unit{kg\cdot s^{-1}}$，$V$为速度，$L=1\unit{m}$，{\small $t_2=1.4\,t_1$}\ $t_{PQ}$过磁场、$MN$过磁场。
% “(未知)”为写在 $V_0$ 下方的小字；$t_2=1.4t_1$ 为写在“$t_{PQ}$过磁场”上方的小字
\par\smallskip
$V_2=\struck{\dfrac{4}{3}}\,\dfrac{2}{3}V_0$（离开磁场时速度）。
% 原稿 $V_2$ 后的 $\frac43$ 疑被涂改/划去
\begin{solution}
\textbf{由平衡条件}
\[
\frac{27}{V_0}=\frac{V_0}{3}\qquad V_0=9\unit{m/s}
\]
\hfill {\small 由于减速了合力向左}\par % 小字，其中“速”字下有下划线
\[
F_A+F'-F=ma\quad\text{以左为正方向}
\]
\[
a=\frac{B^2L^2V}{mR}+\frac{kV}{m}-\frac{P}{mV}
\]
$a$与$v$反向\ 应减速运动\ $a$随之减小，故导体框$PQ$边在磁场% 此行至此为止，原稿未写完

\[
F_A\propto V\qquad F'\propto V
\]
故\ $\dfrac{W_{F'}}{W_A}=\dfrac{F'}{F_A}=\dfrac{kR}{B^2L^2}=4$

由动能定理\quad $W_F-W_{F'}-W_A=\dfrac{1}{2}m\left(\dfrac{2}{3}V_0\right)^2-\dfrac{1}{2}mV_0^2$
\[
W_F=P\cdot 1.4\,\frac{L}{V_0}
\]
\[
W_A=2.64\unit{J}
\]
导体框产生的焦耳热\ $Q=W_A=2.64\unit{J}$
% 页面右侧中下部（y约0.87-0.90）有两道波浪形涂画线，非文字，未录
\end{solution}
\end{problem}
%%%END
%%%BLOCK id=198-01 ch=12 sec=电磁感应·单位长度电阻类 kind=method alt=none cont=none y=0.07-0.37
1、单位长度电阻类

\notefig[0.45]{figures/p198_a.png}

\[
I=\frac{B\,vt\tan\theta\ v}{\dfrac{\rho\left(vt+vt\tan\theta+\dfrac{vt}{\cos\theta}\right)}{S}}
=\frac{BSv\tan\theta}{\rho\left(1+\tan\theta+\dfrac{1}{\cos\theta}\right)}
\]
\[
=\frac{BSv\sin\theta}{\rho(\cos\theta+\sin\theta+1)}\qquad\text{为定值}
\]
%%%END
%%%BLOCK id=198-02 ch=12 sec=电磁感应·单位长度电阻类 kind=method alt=06 cont=none y=0.37-0.82
\notefig[0.4]{figures/p198_b.png}

\[
F=\underset{\text{定值}}{Bi}\ \underset{\uparrow\ \propto t}{L}=Bi\,vt\tan\theta\ \propto t
\]

\notefig[0.4]{figures/p198_c.png}

\struck{$P=P$}\qquad $P=\underline{F}\,v\propto F\propto t$

\notefig[0.4]{figures/p198_d.png}
% 此图纵轴标签为 $P$ 加一个形似小写 n 的下标（疑为 $P_{\text{外}}$ 之类的简写），已含在图内
%%%END
%%%BLOCK id=198-03 ch=12 sec=电磁感应·反向磁场区域线框 kind=note alt=none cont=none y=0.82-0.99
\notefig[0.4]{figures/p198_e.png}

$\uparrow\ \nearrow$\ 2倍电流\ 4倍力（$4F_{A0}$）
% 两个箭头分别指向图中虚线小方框（线框）的两条边
%%%END
%%%BLOCK id=199-01 ch=12 sec=电磁感应·安培力F=kv分类讨论 kind=method alt=03 cont=none y=0.00-0.20
分类讨论\quad $kv$与$F_{\text{安}}$

\notefig[0.28]{figures/p199_a.png}

$F=kv$

\notefig[0.35]{figures/p199_b.png}

\notefig[0.6]{figures/p199_c.png}
% 图中标注：纵轴 $v$、横轴 $t$；三条曲线旁依次为 $k>\frac{B^2L^2}{R}$、$K=\frac{B^2L^2}{R}$（分母 $R$ 被涂改）、$k<\frac{B^2L^2}{R}$

$a=kv$\qquad $a=0$\qquad $v$\unclear{也为}$0$。
%%%END
%%%BLOCK id=199-02 ch=12 sec=电磁感应·i-t图像面积与电荷量 kind=note alt=none cont=none y=0.20-0.24
$i$-$t$图面积\ 若$\Delta\phi=0$\ 则面积$q$应为$0$。
%%%END
